\section{AMC 10 10A 2023}
        \begin{enumerate}
        	\item (\textit{AMC 10 10A 2023}) Cities $A$ and $B$ are $45$ miles apart. Alicia lives in $A$ and Beth lives in $B$. Alicia bikes towards $B$ at $18$ miles per hour. Leaving at the same time, Beth bikes toward $A$ at $12$ miles per hour. How many miles from City $A$ will they be when they meet?


 $\textbf{(A) }20\qquad\textbf{(B) }24\qquad\textbf{(C) }25\qquad\textbf{(D) }26\qquad\textbf{(E) }27$

 \textbf{Answer: }E



\textbf{Solution 1}

This is a $d=st$ problem, so let $x$ be the time it takes to meet. We can write the following equation:
 \[12x+18x=45\]Solving gives us $x=1.5$. The $18x$ is Alicia so $18\times1.5=\boxed{\textbf{(E) 27}}$




\textbf{Solution 2}

The relative speed of the two is $18+12=30$, so $\frac{3}{2}$ hours would be required to travel $45$ miles. $d=st$, so $x=18\cdot\frac{3}{2}=\boxed{\textbf{(E) 27}}$




\textbf{Solution 3}

Since $18$ mph is $\frac{3}{2}$ times $12$ mph, Alicia will travel $\frac{3}{2}$ times as far as Beth. If $x$ is the distance Beth travels, 
 \[\frac{3}{2}x+x=45\] 
 \[\frac{5}{2}x=45\] 
 \[x=18\]Since this is the amount Beth traveled, the amount that Alicia traveled was 
 \[45-18=\boxed{\textbf{(E) 27}}\]




\textbf{Solution 4}

Alice and Barbara close in on each other at 30mph. Since they are 45 miles apart, they will meet in t = d/s = 45miles / 30mph = 3/2 hours. We can either calculate the distance Alice travels at 18mph or the distance Barbara travels at 12mph; since we want the distance from Alice, we go with the former. Alice (and Barbara) will meet in 1 1/2 hours at 18mph x 3/2 hours = 27 miles from A. $\boxed{\textbf{(E) 27}}$




\textbf{Solution 5 (Under 20 seconds)}

We know that Alice approaches Beth at $18$ mph and Beth approaches Alice at $12$ mph. If we consider that if Alice moves $18$ miles at the same time Beth moves $12$ miles —> $15$ miles left. Alice then moves $9$ more miles at the same time as Beth moves $6$ more miles. Alice has moved $27$ miles from point A at the same time that Beth has moved $18$ miles from point B, meaning that Alice and Beth meet $27$ miles from point A.

 



\textbf{Solution 6 (simple linear equations)}

We know that Beth starts 45 miles away from City A, let’s create two equations:Alice-> $18t=d$Beth->  $-12t+45=d$ [-12 is the slope; 45 is the y-intercept]

 Solve the system:

 $18t=-12t+45 30t=45 t=1.5$

 So, $18(1.5)=$ $\boxed{\textbf{(E) 27}}$




\textbf{Solution 7}

Since Alicia and Beth's speeds are constant, they are directly proportional to their distances covered, so the ratio of their speeds is equal to the ratio of their covered distances. Since Alicia travels $\frac{18}{30} = \frac{3}{5}$ of their combined speed, she travels $\frac{3}{5}\cdot 45 = \boxed{\textbf{(E)}\ 27}$ of the total distance.

 -Benedict T (countmath1)

 



\textbf{Solution 8}


                    \begin{figure}[H]
                        \centering
                        \includegraphics[width=0.5\textwidth]{Solution/AMC_10_10A_2023_1_Solution_Figure_1.png}
                    \end{figure}
                    
                    \begin{figure}[H]
                        \centering
                        \includegraphics[width=0.5\textwidth]{Solution/AMC_10_10A_2023_1_Solution_Figure_2.png}
                    \end{figure}
                    Note that Alicia and Beth must have travelled 45 miles together in order for them to meet together. While travelling the 45 miles, Alicia came closer towards Beth and Beth came closer towards Alicia. Since Alicia is travelling faster than Beth, we know that they will meet slightly closer towards city B from the middle. Also, the distance remaining between Alicia and Beth as they bike toward each other is proportional to their combined velocity. The combined velocity is $18t + 12t = 30t$. The time it takes to travel 45 miles going at 30 miles per hour can be calculated using simple algebra.

 Let $t$ be the time it takes for Alicia and Beth to meet together.

 \begin{align*}18t + 12t = 45 \\30t = 45 \\t = \frac{3}{2}\end{align*}


 It takes $\frac{3}{2}$  hours of cycling by both Alicia and Beth at their respective biking speeds in order to meet together. Also, note that the question asks for the distance from City A they will be when they meet. This is the same as the distance travelled by Alicia from City A. Alicia travels at 18 mph, which turns into a distance = speed x time question. 

 \begin{align*}d &= st \\d &= 18 \times \frac{3}{2} \\d &= 27\end{align*}
The distance from City A when Alicia and Beth meet is $\boxed{\textbf{(E)}\ 27}$.

 Note - Even though it might seem like there is a lot of work for this problem, it would typically take a very short time to think this entire process out. 

 


	\item (\textit{AMC 10 10A 2023}) The weight of $\frac{1}{3}$ of a large pizza together with $3 \frac{1}{2}$ cups of orange slices is the same weight of $\frac{3}{4}$ of a large pizza together with $\frac{1}{2}$ cups of orange slices. A cup of orange slices weigh $\frac{1}{4}$ of a pound. What is the weight, in pounds, of a large pizza?


 $\textbf{(A) }1\frac{4}{5}\qquad\textbf{(B) }2\qquad\textbf{(C) }2\frac{2}{5}\qquad\textbf{(D) }3\qquad\textbf{(E) }3\frac{3}{5}$

 \textbf{Answer: }A



\textbf{Solution 1 (Substitution)}

Use a system of equations. Let $x$ be the weight of a pizza and $y$ be the weight of a cup of orange slices.We have 
 \[\frac{1}{3}x+\frac{7}{2}y=\frac{3}{4}x+\frac{1}{2}y.\]Rearranging, we get 
 \begin{align*} \frac{5}{12}x&=3y, \\ x&=\frac{36}{5}y. \end{align*}
Plugging in $\frac{1}{4}$ pounds for $y$ by the given gives $\frac{9}{5}=\boxed{\textbf{(A) }1\frac{4}{5}}.$




\textbf{Solution 2}

Let:$p$ be the weight of a pizza.$o$ be the weight of a cup of orange.

 From the problem, we know that $o = \frac{1}{4}$.

 Write the equation below:
 \[\frac{1}{3} p + \frac{7}{2}\cdot\frac{1}{4} = \frac{3}{4} p + \frac{1}{2}\cdot\frac{1}{4}\]

 Solving for $p$:$\frac{5}{12} p = \frac{3}{4}$
$p = \frac{9}{5} = \boxed{\textbf{(A) }1\frac{4}{5}}.$




\textbf{Solution 3}

$\frac{P}{3} + \frac{7}{2} R = \frac{3}{4} P + \frac{R}{2}$ where $P$ is the pizza weight and $R$ is the weight of cup of orangesSince oranges weigh $\frac{1}{4}$ pound per cup, the oranges on the LHS weigh $\frac{7}{2}$ cups x $\frac{1}{4}$ pounds/cup = $\frac{7}{8}$ pound,and those on the RHS weigh $\frac{1}{2}$ cup x $\frac{1}{4}$ pounds/cup = $\frac{1}{8}$ pound.

 So $\frac{P}{3}$ + $\frac{7}{8}$ pound = $\frac{3}{4} P$ + $\frac{1}{8}$ pound; $\frac{P}{3}$ + $\frac{3}{4}$ pound = $\frac{3}{4} P$.

 Multiplying both sides by $\text{lcm}(3,4) = 12$, we have$4P + 9 = 9P$; $5P = 9$; $P$ = weight of a large pizza = $\frac{9}{5}$ pounds = $\boxed{\textbf{(A)}1 \frac{4}{5}}$ pounds. 

 


	\item (\textit{AMC 10 10A 2023}) How many positive perfect squares less than $2023$ are divisible by $5$? 


 $\textbf{(A) }8\qquad\textbf{(B) }9\qquad\textbf{(C) }10\qquad\textbf{(D) }11\qquad\textbf{(E) }12$

 \textbf{Answer: }A



\textbf{Solution 1 (slightly refined)}

Since $\left \lfloor{\sqrt{2023}}\right \rfloor = 44$, there are $\left \lfloor{\frac{44}{5}}\right \rfloor = \boxed{\textbf{(A) 8}}$ perfect squares less than 2023 that are divisible by 5.

 



\textbf{Solution 2}

Since $5$ is square-free, each solution must be divisible by $5^2=25$. We take $\left \lfloor{\frac{2023}{25}}\right \rfloor = 80$ and see that there are $\boxed{\textbf{(A) 8}}$ positive perfect squares no greater than $80$.

 



\textbf{Solution 3}

Since the perfect squares have to be divisible by 5, then we know it has to be 5 times some number squared $(5 \ cdot x)^{2}$. With this information, you can figure out every single product of 5 and another number squared to count how many perfect squares are divisible by 5 that are less than 2023. (EX: $5^{2} = 25, 10^{2} = 100, ... 40^{2} = 1600$) With this you get a max of $40^{2}$, or  $\left \lfloor{\frac{44}{5}}\right \rfloor = \boxed{\textbf{(A) 8}}$ solutions. 

 



\textbf{Solution 4}

The way of BlueShardow refined:

 All it takes is to recall that 45 squared is 2025, and 45 is 5 x 9. So all the squares of 5 x 1, 5 x 2, 5 x 3 so on are divisible by 5. So the answer is 8. It can be done even if one does not remember that 45 squared is 2025, all it takes is intuition. One can easily see mentally that 5 x 8 that is 40 squared is 1600, and then one has to do just one more computation and see that 5 x 9 that is 45 squared exceeds 2023, so the answer is 8. 

 Note that you can find the square of any number that ends in 5 by taking the number 5 more than it and the number 5 less than it, multiplying those together, and adding 25. For example, to calculate the square of 45, you do 40 x 50 = 2000, and 2000 + 25 = 2025.

 And then you know that it goes up from $1$ to $44$, and the multiples of five are {$5$, $10$..., $35$, $40$}, dividing each of them by $5$, you get {$1$, $2$..., $7$, $8$}, therefore, the answer is $\boxed{\textbf{(A) 8}}$



	\item (\textit{AMC 10 10A 2023}) A quadrilateral has all integer side lengths, a perimeter of $26$, and one side of length $4$. What is the greatest possible length of one side of this quadrilateral?


 $\textbf{(A) }9\qquad\textbf{(B) }10\qquad\textbf{(C) }11\qquad\textbf{(D) }12\qquad\textbf{(E) }13$

 \textbf{Answer: }D



\textbf{Solution 1}

Let's use the triangle inequality. We know that for a triangle, the sum of the 2 shorter sides must always be longer than the longest side. This is because if the longest side were to be as long as the sum of the other sides, or longer, we would only have a line.

 Similarly, for a convex quadrilateral, the sum of the shortest 3 sides must always be longer than the longest side. Thus, the answer is $\frac{26}{2}-1=13-1=\boxed {\textbf{(D) 12}}$

 Sidenote: If there weren't a restriction on integer side lengths, the answer would be the decimal just less than 13, so the sum of the other 3 sides could be just more than 13. That would make the longest side 12.99999..., stopping at who knows how many 9's.

 Sidenote \#2: In the previous sidenote above, it states that the longest side would be 12.999999...., but that is conceptually 13, therefore the answer would be unknown.

 



\textbf{Solution 2}

Say the chosen side is $a$ and the other sides are $b,c,d$.

 By the Generalised Polygon Inequality, $a<b+c+d$. We also have $a+b+c+d=26\Rightarrow b+c+d=26-a$.

 Combining these two, we get $a<26-a\Rightarrow a<13$.

 The largest length that satisfies this is $a=\boxed {\textbf{(D) 12}}$




\textbf{Solution 3}

This is an AMC 10 problem 4, so there is no need for any complex formulas. The largest singular side length from a quadrilateral comes from a trapezoid. So we can set the $2$ sides of the trapezoid equal to $4$. Next we can split the trapezoid into $5$ triangles, where each base length of the triangle equals $4$. So the top side equals $8$, and the bottom side length equals $4+4+4$ $=$ $\boxed {\textbf{(D) 12}}$~ kabbybear

 
                    \begin{figure}[H]
                        \centering
                        \includegraphics[width=0.5\textwidth]{Solution/AMC_10_10A_2023_4_Solution_Figure_1.png}
                    \end{figure}
                    



	\item (\textit{AMC 10 10A 2023}) How many digits are in the base-ten representation of $8^5 \cdot 5^{10} \cdot 15^5$?


 $\textbf{(A) } 14 \qquad\textbf{(B) }15 \qquad\textbf{(C) }16 \qquad\textbf{(D) }17 \qquad\textbf{(E) } 18$

 \textbf{Answer: }E



\textbf{Solution 1}

Prime factorizing this gives us $2^{15}\cdot3^{5}\cdot5^{15}=10^{15}\cdot3^5=243\cdot10^{15}$.

 $10^{15}$ has $16$ digits and $243$ = $2.43*10^{2}$ gives us $3$ more digits. $16+2=\text{\boxed{\textbf{(E) }18}}$
$2.43*10^{17}$ has $18$ digits

 



\textbf{Solution 2 (Only if you don't know how to do the rest of the problems and have about 20 minutes left, IT IS NOT recommended)}

Multiplying it out, we get that $8^5 \cdot 5^{10} \cdot 15^5 = 243000000000000000$. Counting, we have the answer is $\text{\boxed{\textbf{(E) }18}}$~andliu766

 



\textbf{Solution 3 (Similar to Solution 1)}

All the exponents have a common factor of $5$ which we can factor out. This leaves us with $(8 \cdot 5^2 \cdot 15)^5 = (3000)^5 = (3 \cdot 1000)^5$. We can then distribute the power leaving us with $3^5 \cdot 10^{3 \cdot 5} = 243 \cdot 10^{15}$. This would be $243$ followed by $15$ zeros resulting in our answer being $15+3=\text{\boxed{\textbf{(E)}18}}$~leon\_0iler

 



\textbf{Solution 4 (log's)}

We have 8\textasciicircum{}5, 5\textasciicircum{}10, and 15\textasciicircum{}5, if we do a logarithm base of 10 we can figure out how many digits are in the expression when multiplied. So we have $( (log (5^10 \cdot 8^5 \cdot 15^5))$ which using some rules we have $((10log(5))+(5log(8))+(5log(15)))$, using the fact that $log(5)=0.7, log(8)=0.9, log(15)=1.2,$ multiplying it out we get $10*0.7+5*0.9+1.2*5$, and since log(1)=0, as this is false as 1 is a digit and not 0, we add one and then take the floor to get us $\text{\boxed{\textbf{(E)}18}}$~~ Randomuser456All values stated are approximated

 


	\item (\textit{AMC 10 10A 2023}) An integer is assigned to each vertex of a cube. The value of an edge is defined to be the sum of the values of the two vertices it touches, and the value of a face is defined to be the sum of the values of the four edges surrounding it. The value of the cube is defined as the sum of the values of its six faces. Suppose the sum of the integers assigned to the vertices is $21$. What is the value of the cube? 


 $\textbf{(A) } 42 \qquad \textbf{(B) } 63 \qquad \textbf{(C) } 84 \qquad \textbf{(D) } 126 \qquad \textbf{(E) } 252$

 \textbf{Answer: }D



\textbf{Solution 1}

Each of the vertices is counted $3$ times because each vertex is shared by three different edges. Each of the edges is counted $2$ times because each edge is shared by two different faces. Since the sum of the integers assigned to all vertices is $21$, the final answer is $21\times3\times2=\boxed{\textbf{(D) } 126}$




\textbf{Solution 2}

Note that each vertex is counted $2\times 3=6$ times. Thus, the answer is $21\times6=\boxed{\textbf{(D) } 126}$.

 



\textbf{Solution 3}

Just set one vertex equal to $21$, it is trivial to see that there are $3$ faces with value $42$, and $42 \cdot 3=\boxed{\textbf{(D) } 126}$.

 



\textbf{Solution 4}

Since there are 8 vertices in a cube, there are $\dfrac{21}4$ vertices for two edges. There are $4$ edges per face, and $6$ faces in a cube, so the value of the cube is $\dfrac{21}4 \cdot 24 = \boxed{\textbf{(D) } 126}$. 

 (Minor formatting by Technodoggo)

 



\textbf{Solution 5 (use an example)}

Set each vertex to value 1, so the sum of the vertices is 8. We find that the value of the cube, if all vertices are 1, is 48. We conclude that the value of the cube is 6 times the value of the sum of the vertices. Therefore, we choose $21\times6=\boxed{\textbf{(D) } 126}$




\textbf{Solution 6}

The wording of the problem implies that the answer should hold for any valid combination of integers. Thus, we choose the numbers $21, 0, 0, 0, 0, 0, 0, 0$, which are indeed $8$ integers that add to $21$. Doing this, we find three edges that have a value of $21$, and from there, we get three faces with a value of $42$ (while the other three faces have a value of $0$). Adding the three faces together, we get $42+42+42 = \boxed{\textbf{(D) } 126}$.

 



\textbf{Solution 7}

Let the vertices be $a, b, c, d, e, f, g, h$. We know that $a+b+c+d+e+f+g+h=21$. All the edges in a cube are equal, so we can assign the value $x$ to the edge length. 2 vertices make an edge, so $4x=21$. This is also the same value of a face, so we can multiply $21$ by $6$ and get $\boxed{\textbf{(D) } 126}$.

 


	\item (\textit{AMC 10 10A 2023}) Janet rolls a standard $6$-sided die $4$ times and keeps a running total of the numbers she rolls. What is the probability that at some point, her running total will equal $3?$


 $\textbf{(A) }\frac{2}{9}\qquad\jtextbf{(B) }\frac{49}{216}\qquad\textbf{(C) }\frac{25}{108}\qquad\textbf{(D) }\frac{17}{72}\qquad\textbf{(E) }\frac{13}{54}$ (Error compiling LaTeX. Unknown error\_msg)


 
 \textbf{Answer: }B



\textbf{Solution 1 (Casework)}

There are $3$ cases where the running total will equal $3$: one roll; two rolls; or three rolls:

 \textbf{Case 1:}The chance of rolling a running total of $3$, namely $(3)$ in exactly one roll is $\frac{1}{6}$.

 \textbf{Case 2:}The chance of rolling a running total of $3$ in exactly two rolls, namely $(1, 2)$ and $(2, 1)$ is $\frac{1}{6}\cdot\frac{1}{6}\cdot2=\frac{1}{18}$.

 \textbf{Case 3:}The chance of rolling a running total of 3 in exactly three rolls, namely $(1, 1, 1)$ is $\frac{1}{6}\cdot\frac{1}{6}\cdot\frac{1}{6}=\frac{1}{216}$.

 Using the rule of sum we have $\frac{1}{6}+\frac{1}{18}+\frac{1}{216}=\boxed{\textbf{(B) }\frac{49}{216}}$.

 



\textbf{Solution 2 (Brute Force)}

Because there is only a maximum of 3 rolls we must count (running total = 3 means there can't be a fourth roll counted), we can simply list out all of the probabilities.

 If we roll a 1 on the first, the rolls that follow must be 2 or {1,1}, with the following results not mattering. This leaves a probability of $\frac{1}{6}\times\frac{1}{6} + \frac{1}{6}\times\frac{1}{6}\times\frac{1}{6} = \frac{7}{216}$.

 If we roll a 2 on the first, the roll that follows must be 1, resulting in a probability of $\frac{1}{6}\times\frac{1}{6} = \frac{1}{36}$.

 If we roll a 3 on the first, the following rolls do not matter, resulting in a probability of $\frac{1}{6}$.Any roll greater than three will result in a running total greater than 3 no matter what, so those cases can be ignored.Summing the answers, we have $\frac{7}{216} + \frac{1}{36} + \frac{1}{6} = \frac{7+6+36}{216} = \boxed{\textbf{(B) }\frac{49}{216}}$.

 



\textbf{Solution 3}

Consider sequences of 4 integers with each integer between 1 and 6, the number of permutations of 6 numbers is $6^4=1296$. 

 The following 4 types of sequences that might generate a running total of the numbers to be equal to 3 (x, y, or z denotes any integer between 1 and 6).

 Sequence \#1, (1, 1, 1, x): there are $6$ possible sequences. 

 Sequence \#2, (1, 2, x, y): there are $6^2 = 36$ possible sequences. 

 Sequence \#3, (2, 1, x, y): there are $6^2 = 36$ possible sequences. 

 Sequence \#4, (3, x, y, z): there are $6^3 = 216$ possible sequences.

 Out of 1296 possible sequences, there are a total of $6 + 36 + 36 + 216 = 294$ sequences that qualify. Hence, the probability is $294 / 1296 = \boxed{\textbf{(B) }\frac{49}{216}}$




\textbf{Solution 4 (cheese-ish)}

You have to add 1/216 for the case 1-1-1, so it should have denominator 216. Therefore, the answer is B: 49/216

 -Supernova77

 


	\item (\textit{AMC 10 10A 2023}) Barb the baker has developed a new temperature scale for her bakery called the Breadus scale, which is a linear function of the Fahrenheit scale. Bread rises at $110$ degrees Fahrenheit, which is $0$ degrees on the Breadus scale. Bread is baked at $350$ degrees Fahrenheit, which is $100$ degrees on the Breadus scale. Bread is done when its internal temperature is $200$ degrees Fahrenheit. What is this, in degrees, on the Breadus scale?


 $\textbf{(A) }33\qquad\textbf{(B) }34.5\qquad\textbf{(C) }36\qquad\textbf{(D) }37.5\qquad\textbf{(E) }39$

 \textbf{Answer: }D



\textbf{Solution 1 (Substitution)}

To solve this question, you can use $f(x) = mx + b$ where the $x$ is Fahrenheit and the $y$ is Breadus. We have $(110,0)$ and $(350,100)$. We want to find the value of $y$ in $(200,y)$ that falls on this line. The slope for these two points is $\frac{5}{12}$; $y = \frac{5}{12}x + b$. Solving for $b$ using $(110, 0)$, $\frac{550}{12} = -b$. We get $b = \frac{-275}{6}$. Plugging in $(200, y), \frac{1000}{12}-\frac{550}{12}=y$. Simplifying, $\frac{450}{12} = \boxed{\textbf{(D) }37.5}$




\textbf{Solution 2 (Faster)}

Let $^\circ B$ denote degrees Breadus. We notice that $200^\circ F$ is $90^\circ F$ degrees more than $0^\circ B$, and $150^\circ F$ less than $100^\circ B$. This ratio is $90:150=3:5$; therefore, $200^\circ F$ will be $\dfrac3{3+5}=\dfrac38$ of the way from $0$ to $100$, which is $\boxed{\textbf{(D) }37.5}$




\textbf{Solution 3 (Intuitive)}

From $110$ to $350$ degrees Fahrenheit, the Breadus scale goes from $1$ to $100$. $110$ to $350$ degrees is a span of $240$, and we can use this to determine how many Fahrenheit each Breadus unit is worth. $240$ divided by $100$ is $2.4$, so each Breadus unit is $2.4$ Fahrenheit, starting at $110$ Fahrenheit. For example, $1$ degree on the Breadus scale is $110 + 2.4$, or $112.4$ Fahrenheit. Using this information, we can figure out how many Breadus degrees $200$ Fahrenheit is. $200-110$ is $90$, so we divide $90$ by $2.4$ to find the answer, which is $\boxed{\textbf{(D) }37.5}$




\textbf{Solution 4}

We note that the range of F temperatures that $0-100$ $\text{Br}^\circ$ represents is $350-110 = 240$ $\text{F}^\circ$.$200$ $\text{F}^\circ$ is $(200-110) = 90$ $\text{F}^\circ$ along the way to getting to $240$ $\text{F}^\circ$, the end of this range, or $90/240 = 9/24 = 3/8 = 0.375$ of the way. Therefore if we switch to the Br scale, we are $0.375$ of the way to $100$ from $0$, or at $\boxed{\textbf{(D) }37.5}$ $\text{Br}^\circ$.

 



\textbf{Solution 5}

We have the points $(0, 110)$ and $(100, 350)$. We want to find $(x, 200)$. The equation of the line is $y=\frac{12}{5}x+110$. We use this to find $x=\frac{75}{2}=37.5$, or $\boxed{D}$.~MC413551

 



\textbf{Solution 6 (extremely simple)}

We can write the value $y$ on the Breadus scale as $y = mt + b$, where $t$ is the temperature in Fahrenheit.From the problem, $110m + 1b = 0$ and $350m + 1b = 100.$We can rewrite this problem in terms of linear algebra to solve it.

 $Let \: A =\begin{bmatrix} 110 & 1 \\ 350 & 1 \end{bmatrix}, let \: B = \begin{bmatrix} 0 \\ 100 \end{bmatrix}, and \: let \: x = \begin{bmatrix} m \\ b \end{bmatrix}.$We can write the system of equations as Ax = B.We can solve for x using the expression x = $A^{-1}B$.Calculating this value we get $x = \begin{bmatrix} -1/240 & 1/240 \\ 35/24 & -11/24 \end{bmatrix}\cdot\begin{bmatrix} 0 \\ 100 \end{bmatrix}=\begin{bmatrix} 5/12 \\ -275/6 \end{bmatrix}.$Therefore, $m = 5/12 \: and \: b = -275/6$.Plugging in $t = 200$, we get $(5/12)200+(-275/6) = \boxed{\textbf{(D) }37.5}$.

 


	\item (\textit{AMC 10 10A 2023}) A digital display shows the current date as an $8$-digit integer consisting of a $4$-digit year, followed by a $2$-digit month, followed by a $2$-digit date within the month. For example, Arbor Day this year is displayed as $20230428.$ For how many dates in $2023$ does each digit appear an even number of times in the $8$-digital display for that date?


 $\textbf{(A)}~5\qquad\textbf{(B)}~6\qquad\textbf{(C)}~7\qquad\textbf{(D)}~8\qquad\textbf{(E)}~9$

 \textbf{Answer: }E



\textbf{Solution 1 (Casework)}

Do careful casework by each month. Make sure to start with 2023. In the month and the date, we need a $0$, a $3$, and two digits repeated (which has to be $1$ and $2$ after consideration). After the casework, we get $\boxed{\textbf{(E)}~9}$. For curious readers, the numbers (in chronological order) are:

 \begin{verbatim}
202301132023013120230223202303112023032220231013202310312023110320231130
\end{verbatim}





\textbf{Solution 2}

There is one $3$, so we need one more (three more means that either the month or units digit of the day is $3$). For the same reason, we need one more $0$. By inspection, we can see that the other two digits must be a repeated digit that is not $3$ or $0$. 

 
If $3$ is the units digit of the month, then the $0$ must be the units digit of the month, since there is no month ending with 3 that is more than March. In this case, the last two digits must match ($11, 22$). This gives $2023-03-11$ and $2023-03-22$, so 2 successes.

 If $3$ is the tens digit of the day, then $0$ can be in any of the remaining three places. If it is in the ones digit of the the day, then the month must be $11$, which gives $2023-11-30$. If it is in the ones digit of the month, then the month must be $10$, in which case the day must be $31$, which gives $2023-10-31$. If $0$ is in the tens digit of the month, then the month must be $01$, which gives $2023-01-31$. Together, $2023-11-30$, $2023-10-31$, and $2023-01-31$ make 3 successes. 

 
If $3$ is the units digit of the day, then $0$ could go in any of the $3$ remaining slots again. If it's the tens digit of the day, then the other digits must be $1$. If $0$ is the units digit of the month, then the other two slots must both be $1$. If $0$ is the tens digit of the month, then the other two slots can be either both $1$ or both $2$. In total, there are $4$ successes here. 

 Summing through all cases, there are $2 + 3 + 4 = \boxed{\textbf{(E)}~9}$ dates. 

 -Benedict T (countmath1)

 -Edits by Cat\_Of\_Logic

 



\textbf{Solution 3}

We start with $2023----$ we need an extra $0$ and an extra $3$. So we have at least one of those extras in the days, except we can have the month $03$. We now have $6$ possible months $01,02,03,10,11,12$. For month $1$ we have two cases, we now have to add in another 1, and the possible days are $13,31$. For month $2$ we need an extra $2$ so we can have the day $23$ note that we can't use $32$ because it is to large. Now for month $3$ we can have any number and multiply it by $11$ so we have the solution $11,22$. For October we need a $1$ and a $3$ so we have $13,31$ as our choices. For November we have two choices which are $03,30$.Now for December we have $0$ options. Summing $2+1+2+2+2$ we get $\boxed{\textbf{(E)}~9}$ solutions.

 


	\item (\textit{AMC 10 10A 2023}) Maureen is keeping track of the mean of her quiz scores this semester. If Maureen scores an $11$ on the next quiz, her mean will increase by $1$. If she scores an $11$ on each of the next three quizzes, her mean will increase by $2$. What is the mean of her quiz scores currently? 


 $\textbf{(A) }4\qquad\textbf{(B) }5\qquad\textbf{(C) }6\qquad\textbf{(D) }7\qquad\textbf{(E) }8$

 \textbf{Answer: }D



\textbf{Solution 1}

Let $a$ represent the amount of tests taken previously and $x$ the mean of the scores taken previously. 

 We can write the following equations:

 
 \[\frac{ax+11}{a+1}=x+1\qquad (1)\]
 \[\frac{ax+33}{a+3}=x+2\qquad (2)\]

 Multiplying equation $(1)$ by $(a+1)$ and solving, we get:
 \[ax+11=ax+a+x+1\]
 \[11=a+x+1\]
 \[a+x=10\qquad (3)\]

 Multiplying equation $(2)$ by $(a+3)$ and solving, we get:
 \[ax+33=ax+2a+3x+6\]
 \[33=2a+3x+6\]
 \[2a+3x=27\qquad (4)\]

 Solving the system of equations for $(3)$ and $(4)$, we find that $a=3$ and $x=\boxed{\textbf{(D) }7}$.

 



\textbf{Solution 2 (Variation on Solution 1)}

Suppose Maureen took $n$ tests with an average of $m$.

 If she takes another test, her new average is $\frac{(nm+11)}{(n+1)}=m+1$

 Cross-multiplying: $nm+11=nm+n+m+1$, so $n+m=10$.

 If she takes $3$ more tests, her new average is $\frac{(nm+33)}{(n+3)}=m+2$

 Cross-multiplying: $nm+33=nm+2n+3m+6$, so $2n+3m=27$.

 But $2n+3m$ can also be written as $2(n+m)+m=20+m$. Therefore $m=27-20=\boxed{\textbf{(D) }7}$




\textbf{Solution 3 (do this if you are bored)}

Let $s$ represent the sum of Maureen's test scores previously and $t$ be the number of scores taken previously.

 So, $\frac{s+11}{t+1} = \frac{s}{t}+1$ and $\frac{s+33}{t+3} = \frac{s}{t}+2$

 We can use the first equation to write $s$ in terms of $t$.

 We then substitute this into the second equation: $\frac{-t^2+10t+33}{t+3} = \frac{-t^2+10t}{t}+2$

 From here, we should solve for t: multiply both sides by ($t$) and then ($t+3$), combining like terms to get $t^2-3t=0$. Factorize to get $t=0$ or $t=3$, and therefore $t=3$ (makes sense for the problem).

 We substitute this to get $s=21$.

 Therefore, the solution to the problem is $\frac{21}{3}=$ $\boxed{\textbf{(D) }7}$




\textbf{Solution 4 (Testing Answer Choices)}

Let's consider all the answer choices. If the average is $8$, then, we can assume that all her test choices were $8$. We can see that she must have gotten $8$ twice, in order for another score of $11$ to bring her average up by one. However, adding three $11$'s will not bring her score up to 10. Continuing this process for the answer choices, we see that the answer is $\boxed{\textbf{(D) }7}$~andliu766 (minor wording edit by mihikamishra)

 



\textbf{Solution 5}

Let $n$ be the number of existing quizzes. So after one more test, score $11$ has $n+1$ extra points to distribute to $n+1$ quizzes. Also, after three more quizzes, there will be $3(n+1)$ extra points to distribute to the $n+3$ quizzes. So $3n+3=2(n+3)$. This means $n=3$. $n+1$ extra points means original mean (average) is $7$



	\item (\textit{AMC 10 10A 2023}) A square of area $2$ is inscribed in a square of area $3$, creating four congruent triangles, as shown below. What is the ratio of the shorter leg to the longer leg in the shaded right triangle? 


 
                    \begin{figure}[H]
                        \centering
                        \includegraphics[width=0.5\textwidth]{Problem/AMC_10_10A_2023_11_Problem_Figure_1.png}
                    \end{figure}
                    
$\textbf{(A) }\frac15\qquad\textbf{(B) }\frac14\qquad\textbf{(C) }2-\sqrt3\qquad\textbf{(D) }\sqrt3-\sqrt2\qquad\textbf{(E) }\sqrt2-1$

 \textbf{Answer: }C



\textbf{Solution 1}

The side lengths of the inner square and outer square are $\sqrt{2}$ and $\sqrt{3}$ respectively. Let the shorter side of our triangle be $x$, thus the longer leg is $\sqrt{3}-x$. Hence, by the Pythagorean Theorem, we have 
 \begin{align*} (\sqrt{3}-x)^2+x^2&=(\sqrt{2})^2 \\ 3-2\sqrt{3}x+x^2+x^2&=2 \\ 2x^2-2\sqrt{3}x+1&=0. \end{align*}
By the quadratic formula, we find that $x=\frac{\sqrt{3}\pm1}{2}$, so the answer is $\frac{\sqrt{3}-1}{\sqrt{3}+1}=\boxed{\textbf{(C) }2-\sqrt3}.$




\textbf{Solution 2 (Area)}

Looking at the diagram, we know that the square inscribed in the square with area $3$ has area $2$. Thus, the area outside of the small square is $3-2=1.$ This area is composed of $4$ congruent triangles, so we know that each triangle has an area of $\dfrac14$. 

 From solution $1$, the base (short side of the triangle) has a length $x$ and the height $\sqrt{3} - x$, which means that $\frac{x\left(\sqrt{3} - x\right)}{2} = \frac{1}{4}$. 

 We can turn this into a quadratic equation: $x^2-x\sqrt{3}+\frac{1}{2} = 0$. 

 By using the quadratic formula, we get $x=\frac{\sqrt{3}\pm1}{2}$.Therefore, the answer is $\frac{\sqrt{3}-1}{\sqrt{3}+1}=\boxed{\textbf{(C) }2-\sqrt3}.$

 (Clarity & formatting edits by Technodoggo)

 



\textbf{Solution 3}

Let $x$ be the ratio of the shorter leg to the longer leg, and $y$ be the length of longer leg. The length of the shorter leg will be $xy$.

 Because the sum of two legs is the side length of the outside square, we have $xy + y = \sqrt{3}$, which means $(xy)^2 + y^2 + 2xy^2 = 3$. Using the Pythagorean Theorem for the shaded right triangle, we also have $(xy)^2 + y^2 = 2$. Solving both equations, we get $2xy^2 = 1$. Using $y^2=\frac{1}{2x}$ to substitute $y$ in the second equation, we get $x^2\cdot \frac{1}{2x} + \frac{1}{2x} = 2$. Hence, $x^2 - 4x + 1 = 0$. By using the quadratic formula, we get $x=2\pm \sqrt3$. Because $x$ be the ratio of the shorter leg to the longer leg, it is always less than $1$. Therefore, the answer is $\boxed{\textbf{(C) }2-\sqrt3}$.

 



\textbf{Solution 4}

The side length of the bigger square is equal to $\sqrt{3}$, while the side length of the smaller square is $\sqrt{2}$. Let $x$ be the shorter leg and $y$ be the longer one. Clearly, $x+y=\sqrt{3}$, and $xy=\frac{1}{2}$. Using Vieta's to build a quadratic, we get 
 \[x^2-\sqrt{3}x+\frac{1}{2}=0.\] Solving, we get $x=\frac{\sqrt{3}-1}{2}$ and $y=\frac{\sqrt{3}+1}{2}$. Thus,  we find $\frac{x}{y}=\frac{\sqrt{3}-1}{\sqrt{3}+1}=\frac{(\sqrt{3}-1)^2}{(\sqrt{3}+1)\cdot(\sqrt{3}-1)}=\frac{4-2\sqrt{3}}{2}=\boxed{\textbf{(C) }2-\sqrt3}$.

 


 



\textbf{Solution 5}

Let $\theta$ be the angle opposite the smaller leg. We want to find $\tan\theta$.

 The area of the triangle is $\frac{1}{2}\left(\sqrt{2}\sin\theta\right)\left(\sqrt{2}\cos\theta\right)=\frac{1}{2}\sin 2\theta=\frac{1}{4},$ which implies $\sin 2\theta = \frac{1}{2},$ or $\theta=15^\circ$. Therefore $\tan \theta = \boxed{\textbf{(C) }2-\sqrt3}$




\textbf{Solution 6}

Allow a and b to be the sides of a triangle. WLOG, suppose $a > b.$We want to find $\frac{b}{a}$. Notice that the area of a triangle is $\frac{3-2}{4}$, which results in $\frac{1}{4}$. Thus, $ab = \frac{1}{2}$. However, the square of the hypotenuse of this triangle is $a^2+b^2$, but also $2$. We can write $b$ as $\frac{1}{2a}$, and then plug it in. We get $a^2+\frac{1}{4a^2} = 2$, so $4a^4-8a^2+1 = 0$. Applying the quadratic formula, $a^2 = \frac{8 \pm 4\sqrt{3}}{2}$, or $4 \pm 2\sqrt3$. However, since $a$ and $b$ must both be solutions of the quadratic, since both equations were cyclic. Since $a>b$, then $a^2 = 4+2\sqrt3$, and $b^2 = 4-2\sqrt3$. To find $\frac{b}{a}$, we can simply find the square root of $\frac{b^2}{a^2}$. This is $\sqrt{\frac{4-2\sqrt3}{4+2\sqrt3}} = \sqrt{\frac{2-\sqrt3}{2+\sqrt3}} = \sqrt{(2-\sqrt3)(2-\sqrt3)} = \boxed {2-\sqrt3}$, so the answer is $\boxed {C}$. - Sepehr2010

 



\textbf{Solution 7 (Manipulation)}

Let $a$ be the length of the shorter leg and $b$ be the longer leg. By the Pythagorean theorem, we can derive that $a^2+b^2=2$. Using area we can also derive that $2ab=1$. $(a+b)^2=3$ as given in the diagram, we can find that $(a-b)^2=1$ because $4ab=2$. This means that $b+a=\sqrt{3}$ and $b-a=1$. Adding the equations gives $b=\frac{\sqrt{3}+1}{2}$ and when $b$ is plugged in $a = \frac{\sqrt{3}-1}{2}$. Rationalizing the denominators gives us $\boxed{\textbf{(C) } 2-\sqrt{3}}$.

 



\textbf{Solution 8}

Let $s$ be the length of the shorter leg, $l$ the longer leg, and $t$ the desired ratio. We have:
 \[s^2+l^2=2 \text{ and } \frac{1}{2}sl=\frac{1}{4}\]The former dividing the latter gives:
 \[t+\frac{1}{t}=4\]Solve this for $t$.Multiply by $t$:
 \[t^2+1=4t\]Subtract -4t to get a quadratic:
 \[t^2-4t+1=0\] Use the quadratic formula to get:
 \[2 \pm \sqrt{3}\]And since the smaller one is the first one, the solution has to be less than 1, so the answer is the one less than 1, the answer is$\boxed{\textbf{(C) }2-\sqrt3}$



	\item (\textit{AMC 10 10A 2023}) How many three-digit positive integers $N$ satisfy the following properties?


 
            \begin{itemize}
        \item The number $N$ is divisible by $7$.
            \end{itemize}
            
            \begin{itemize}
        \item The number formed by reversing the digits of $N$ is divisible by $5$.
            \end{itemize}
            $\textbf{(A) } 13 \qquad \textbf{(B) } 14 \qquad \textbf{(C) } 15 \qquad \textbf{(D) } 16 \qquad \textbf{(E) } 17$

 \textbf{Answer: }B



\textbf{Solution 1}

Multiples of $5$ will always end in $0$ or $5$, and since the numbers have to be a three-digit numbers, it cannot start with 0 (otherwise it would be a two-digit number), narrowing our choices to 3-digit numbers starting with $5$. Since the numbers must be divisible by 7, all possibilities have to be in the range from $7 \cdot 72$ to $7 \cdot 85$ inclusive(504 to 595). 

 (Add 1 to include 72)

 $85 - 72 + 1 = 14$. $\boxed{\textbf{(B) } 14}$.

 You can also take 497 away from each of the numbers(removing the hundreds digit and adding three to each of the numbers), resulting in the numbers {7, 14, 21..., 84, 91, 98}. Dividing each of them by 7, you get the numbers {1, 2, 3..., 12, 13, 14}. Therefore, the answer is $\boxed{\textbf{(B) 14}}$




\textbf{Solution 2 (solution 1 but more thorough)}

Let $N=\overline{cab}=100c+10a+b.$ We know that $\overline{bac}$ is divisible by $5$, so $c$ is either $0$ or $5$. However, since $c$ is the first digit of the three-digit number $N$, it can not be $0$, so therefore, $c=5$. Thus, $N=\overline{5ab}=500+10a+b.$ There are no further restrictions on digits $a$ and $b$ aside from $N$ being divisible by $7$. 

 The smallest possible $N$ is $504$. The next smallest $N$ is $511$, then $518$, and so on, all the way up to $595$. Thus, our set of possible $N$ is $\{504,511,518,\dots,595\}$. Dividing by $7$ for each of the terms will not affect the cardinality of this set, so we do so and get $\{72,73,74,\dots,85\}$. We subtract $71$ from each of the terms, again leaving the cardinality unchanged. We end up with $\{1,2,3,\cdots,14\}$, which has a cardinality of $14$. Therefore, our answer is $\boxed{\textbf{(B) } 14.}$




\textbf{Solution 3 (modular arithmetic)}

We first proceed as in the above solution, up to $N=500+10a+b$. We then use modular arithmetic: 

 
 \begin{align*} 0&\equiv N \:(\text{mod }7)\\ &\equiv500+10a+b\:(\text{mod }7)\\ &\equiv3+3a+b\:(\text{mod }7)\\ 3a+b&\equiv-3\:(\text{mod }7)\\ &\equiv4\:(\text{mod }7)\\ \end{align*}


 We know that $0\le a,b<10$. We then look at each possible value of $a$: 

 If $a=0$, then $b$ must be $4$. 

 If $a=1$, then $b$ must be $1$ or $8$. 

 If $a=2$, then $b$ must be $5$. 

 If $a=3$, then $b$ must be $2$ or $9$. 

 If $a=4$, then $b$ must be $6$. 

 If $a=5$, then $b$ must be $3$. 

 If $a=6$, then $b$ must be $0$ or $7$. 

 If $a=7$, then $b$ must be $4$. 

 If $a=8$, then $b$ must be $1$ or $8$. 

 If $a=9$, then $b$ must be $5$. 

 Each of these cases are unique, so there are a total of $1+2+1+2+1+1+2+1+2+1=\boxed{\textbf{(B) } 14.}$




\textbf{Solution 4}

The key point is that when reversed, the number must start with a $0$ or a $5$ based on the second restriction. But numbers can't start with a $0$.

 So the problem is simply counting the number of multiples of $7$ in the $500$s.

 $7 \times 72 = 504$, so the first multiple is $7 \times 72$.

 $7 \times 85 = 595$, so the last multiple is $7 \times 85$.

 Now, we just have to count $7\times 72, 7\times 73, 7\times 74,\cdots, 7\times 85$.

 We have a set that numbers $85-71=\boxed{\textbf{(B) 14}}$



	\item (\textit{AMC 10 10A 2023}) Abdul and Chiang are standing $48$ feet apart in a field. Bharat is standing in the same field as far from Abdul as possible so that the angle formed by his lines of sight to Abdul and Chiang measures $60^\circ$. What is the square of the distance (in feet) between Abdul and Bharat?


 $\textbf{(A) } 1728 \qquad \textbf{(B) } 2601 \qquad \textbf{(C) } 3072 \qquad \textbf{(D) } 4608 \qquad \textbf{(E) } 6912$

 \textbf{Answer: }C



\textbf{Solution 1 (Using Trigonometry)}


                    \begin{figure}[H]
                        \centering
                        \includegraphics[width=0.5\textwidth]{Solution/AMC_10_10A_2023_13_Solution_Figure_1.png}
                    \end{figure}
                    

 Let $\theta=\angle ACB$ and $x=\overline{AB}$. 

 According to the Law of Sines, we know that $\dfrac{\sin\theta}x=\dfrac{\sin60^\circ}{48}=\dfrac{\sqrt3}{96}$. Rearranging, we get that $x=\dfrac{\sin\theta}{\frac{\sqrt3}{96}}=32\sqrt3\sin\theta$ where $x$ is a function of $\theta$. We want to maximize $x$. 

 We know that the maximum value of $\sin\theta=1$, so this yields $x=32\sqrt3\implies x^2=\boxed{\textbf{(C) }3072.}$ 

 A quick check verifies that $\theta=90^\circ$ indeed works. 

 



\textbf{Solution 2 (Inscribed Angles)}

We can draw a circle such that the chord AC inscribes an arc of 120 degrees. This way, any point B on the circle not in the inscribed arc will form an angle of 60 degrees with $\angle{ABC}$. To maximize the distance between A and B, they must be opposite each other. So, the problem is now finding the length of the diameter of the circle. We know AOC is 120 degrees, so dropping a perpendicular form O to AC gives us the radius as $16\sqrt{3}$. So, the diameter is $32\sqrt{3}$ which gives us the answer $\boxed{\textbf{(C) }3072}$




\textbf{Solution 3 (Guessing)}

Guess that the optimal configuration is a 30-60-90 triangle, as an equilateral triangle gives an answer of $48^2=2304$, which is not on the answer choices. Its ratio is $\frac{48}{\sqrt{3}}$, so $\overline{AB}=\frac{96}{\sqrt{3}}$.

 Its square is then $\frac{96^2}{3}=\boxed{\textbf{(C) }3072}$

 Note: The distance between Abdul and Chiang is constant, so let that be represented as ${x}$. If we were dealing with an equilateral triangle, the height would be ${{x\sqrt3}/2}$, and if we were dealing with a 30-60-90 triangle, the height would be ${x\sqrt3}$, which is greater than ${{x\sqrt3}/2}$.

 Note: The problem states that the distance from A to B is maximized, and the longest side of a triangle will always correspond to the largest possible angle, aka 90 so 90 and 60 makes it a 30-60-90.

 



\textbf{Solution 4}

We use $A$, $B$, $C$ to refer to Abdul, Bharat and Chiang, respectively.We draw a circle that passes through $A$ and $C$ and has the central angle $\angle AOC = 60^\circ \cdot 2$.Thus, $B$ is on this circle.Thus, the longest distance between $A$ and $B$ is the diameter of this circle.Following from the law of sines, the square of this diameter is
 \[ \left( \frac{48}{\sin 60^\circ} \right)^2 = \boxed{\textbf{(C) 3072}}. \]




\textbf{Solution 5}

We can represent Abdul, Bharat and Chiang as $A$, $B$, and $C$, respectively.Since we have $\angle ABC=60^\circ$ and  (from other solutions) $\angle BCA=90^\circ$, this is a $30-60-90$ triangle.By the side ratios of a $30-60-90$ triangle, we can infer that $AB=\frac{48\times 2}{\sqrt{3}}$.Squaring AB we get $\boxed{\textbf{(C) 3072}}$.

 



\textbf{Solution 6 (easier to visualize)}

Let Abdul and Chiang be a line segment, specifically a chord on a circle. Let chord $\overline{AC}$ be $48$. If Bharat is the furthest away from Abdul, chord $\overline{AB}$ should be the diameter of the circle, as by definition the diameter is the longest chord in a circle. We know that $\triangle {ABC}$ is right since by definition a triangle inscribed in a semicircle is right. Chord $\overline{AC}=48$, $\overline{AB}$ is the hypotenuese, and $\angle{B}=60^{\circ}$. Therefore, the correct answer is $(\frac{48\times{2}}{\sqrt{3}})^2=\boxed{(C) 3072}$




\textbf{Solution 7 (Definition of a Circle)}

The definition of a circle directly shows a circle can circumscribe any three points, and those points will always form a triangle.

 Consider the circle $\Omega$ inscribing triangle $ABC$. It follows that definition via the problem angle $ABC = 60^\circ$. Also, $AC = 48$.

 $\textbf{Lemma}$:

 The only way for $B$ to be as far from $A$ as possible with angle $ABC = 60^\circ$ is only plausible if $AB$ is the diameter of $\Omega$ and $C$ is placed on one semicircle above or below $AB$ so that $ABC = 60^\circ$.

 $\textbf{Proof of the Lemma}$:

 It is by definition of a circle the farthest position from two points on a locus that is a circle must be the circle's diameter. Thus, $AB$ is the diameter of the circle. From here, the angle from $ACB$ through the inscribed angles theorem is $90^\circ$, meaning a position does exist with $ABC = 60^\circ$, when $BAC = 30^\circ$, hence proving the lemma.

 $\textbf{2023 AMC 10A Q13}$:

 Here, we provide the same scenario but with $AC = 48$. Through the properties of 30-60-90 degree triangles one finds that $AB = \frac{48}{\sqrt{3}}$. Then, $AB^2 = 96 \cdot 32 = \boxed{3072}$.

 



\textbf{Solution 8 (Power of a Point, Formal Definition)}

Consider the field to be a circle $\Omega$ with centre $O$ and a points $A$ on the circle. We require $B$ to exist on this plane such that there exists a maximal value between A and B. 

 $\textbf{Lemma 1}$:

 For $AB$ to be maximized, it is required that $B$ lies on the circumference of $\Omega$ and is diametrically opposite to $A$.

 $\textbf{Proof of the Lemma 1}$:

 Obviously, let the radius be $OA = r$ and through Power of a Point on $A$ with respect to $\Omega$ we see that $OA^2 - r^2 = 0$ given definition of A, meaning $A$ indeed lies on the circle. Then, Power of a Point on $B$ with respect to $\Omega$ gives $OB^2 - OA^2 \ge 0$ due to the condition that $B$ lies on the field. This implies that the distance $OB > OA$ or $OB = OA$. Given that $OB > OA$ and $OB^2 - OA^2 > 0$ it is defined that $OA+OB < 2r$. On the contrary, given that $OB = OA$, we see that $OB^2 - OA^2 = 0$ and then $OA + OB = 2r$, which is maximal as through definition of a circle, $2r$ is the maximal plausible distance between two points.

 $\textbf{Lemma 2}$:

 $C$ must lie on $\Omega$
$\textbf{Proof of the Lemma 2}$:

 Say the contrary. It is given that $C$ must lie somewhere on the field. Then, Power of a point on $C$ with respect to $\Omega$ yields $OC^2 - r^2 \ge 0$. Given Lemma 1 maximality occurs if $A$ and $B$ are diametrically opposite. Then, if $C$ lies in the closed disc of $\Omega$, it follows that $ABC = 60^\circ$ given $AB = 2r$. Then, the only plausible outcome is for $C$ to lie outside the disc outside $\Omega$ to create an equilateral triangle of side length $2r$. This contradicts the idea that $C$ lies on the field, and therefore $C$ cannot exist in the closed disc, and must lie on $\Omega$'s circumference, hence proved. Even so, we have to find if such a configuration exist. Given $C$ lies on $\Omega$'s circumference, through the inscribed angles theorem is $90^\circ$, meaning a position does exist with $ABC = 60^\circ$, when $BAC = 30^\circ$, hence showing such a case does exist, proving the lemma.

 $\textbf{2023 AMC 10A Q13}$:

 Here, we provide the same scenario of a combination of lemma 1 and 2 but with $AC = 48$. Through the properties of 30-60-90 degree triangles one finds that $AB = \frac{48}{\sqrt{3}}$. Then, $AB^2 = 96 \cdot 32 = \boxed{3072}$.

 


	\item (\textit{AMC 10 10A 2023}) A number is chosen at random from among the first $100$ positive integers, and a positive integer divisor of that number is then chosen at random. What is the probability that the chosen divisor is divisible by $11$?


 $\textbf{(A)}~\frac{4}{100}\qquad\textbf{(B)}~\frac{9}{200} \qquad \textbf{(C)}~\frac{1}{20} \qquad\textbf{(D)}~\frac{11}{200}\qquad\textbf{(E)}~\frac{3}{50}$

 \textbf{Answer: }B



\textbf{Solution 1}

In order for the divisor chosen to be a multiple of $11$, the original number chosen must also be a multiple of $11$. Among the first $100$ positive integers, there are 9 multiples of 11; 11, 22, 33, 44, 55, 66, 77, 88, 99. We can now perform a little casework on the probability of choosing a divisor which is a multiple of 11 for each of these 9, and see that the probability is 1/2 for each. The probability of choosing these 9 multiples in the first place is $\frac{9}{100}$, so the final probability is $\frac{9}{100} \cdot \frac{1}{2} = \frac{9}{200}$, so the answer is $\boxed{\textbf{(B)}~\frac{9}{200}}.$
$11 = 1, 11 \Rightarrow \frac{1}{2}\\ 22 = 2 \times 11: 1, 2, 11, 22 \Rightarrow \frac{1}{2}\\ 33 = 3 \times 11: 1, 3, 11, 33 \Rightarrow \frac{1}{2}\\ 44 = 2^2 \times 11: 1, 2, 4, 11, 22, 44 \Rightarrow \frac{1}{2}\\ 55 = 5 \times 11: 1, 5, 11, 55 \Rightarrow \frac{1}{2}\\ 66 = 2 \times 3 \times 11: 1, 2, 3, 6, 11, 22, 33, 66 \Rightarrow \frac{1}{2}\\ 77 = 7 \times 11: 1, 7, 11, 77 \Rightarrow \frac{1}{2}\\ 88 = 2^3 \times 11: 1, 2, 4, 8, 11, 22, 44, 88 \Rightarrow \frac{1}{2}\\ 99 = 3^2 \times 11: 1, 3, 9, 11, 33, 99 \Rightarrow \frac{1}{2}$




\textbf{Solution 2}

There are $\left\lfloor\frac{100}{11}\right\rfloor = 9$ multiples of $11$ under $100$. Because all of these numbers are multiples of $11$ to the first power and first power only, their factors can either have $11$ as a factor ($11^{1}$) or not have $11$ as a factor ($11^{0}$), since 11 is prime, resulting in a $\frac{1}{2}$ chance of a factor chosen being divisible by $11$. The chance of choosing any factor of $11$ under $100$ is $\frac{9}{100}$, so the final answer is $\frac{9}{100} \cdot \frac{1}{2} = \boxed{\textbf{(B)}~\frac{9}{200}}.$




\textbf{Solution 3 (generalized)}

Let $N$ be the positive integer in question. Since $N$ is a multiple of $11$, we can write $N = 11^{e_1}p_1^{e_2}\cdots p_{k-1}^{e_k},$ the prime factorization of $N$. To find the total number of divisors divisible by $11$, observe that there are $(e_2 + 1)(e_3 + 1)\cdots (e_k  + 1)$ divisors not divisible by $11$. For each power of $11$ greater than $1$, of which there are $e_1$, it can be paired with any of these other divisors. Therefore, there are 
 \[e_1(e_2+1)(e_3 + 1)\cdots (e_k + 1)\]divisors of $N$ that are divisible by $11$. The total number of divisors of $N$ is 
 \[(e_1 + 1)(e_2 + 1)\cdots (e_k + 1),\] so the probability of choosing a divisor that is divisible by $11$ is 
 \[\frac{e_1(e_2+1)(e_3 + 1)\cdots (e_k + 1)}{(e_1 + 1)(e_2 + 1)\cdots (e_k + 1)} = \frac{e_1}{e_1 + 1}.\]For each positive integer less than or equal to $100$, the highest power of $11$ that divides it is $11$, so $e_1 = 1$, meaning that the probability of choosing a divisor (that is divisible by $11$) of a fixed $N$ is $\frac{1}{2}.$ The probability of choosing any $N$ from the first $100$ positive integers is $\frac{1}{100},$ so the probability of choosing any of these divisors is $\frac{1}{100}\cdot \frac{1}{2}.$ There are $9$ multiples of $11$ less than or equal to $100$, so the total probability is
 \[9\cdot \frac{1}{100}\cdot \frac{1}{2} = \boxed{\textbf{(B)}\ \frac{9}{200}}.\]

 -Benedict T (countmath1)

 


	\item (\textit{AMC 10 10A 2023}) An even number of circles are nested, starting with a radius of $1$ and increasing by $1$ each time, all sharing a common point. The region between every other circle is shaded, starting with the region inside the circle of radius $2$ but outside the circle of radius $1.$ An example showing $8$ circles is displayed below. What is the least number of circles needed to make the total shaded area at least $2023\pi$? 


 
                    \begin{figure}[H]
                        \centering
                        \includegraphics[width=0.5\textwidth]{Problem/AMC_10_10A_2023_15_Problem_Figure_1.png}
                    \end{figure}
                    
$\textbf{(A) } 46 \qquad \textbf{(B) } 48 \qquad \textbf{(C) } 56 \qquad \textbf{(D) } 60 \qquad \textbf{(E) } 64$

 \textbf{Answer: }E



\textbf{Solution 1}

Notice that the area of the shaded region is $(2^2\pi-1^2\pi)+(4^2\pi-3^2\pi)+(6^2\pi-5^2\pi)+ \cdots + (n^2\pi-(n-1)^2 \pi)$ for any even number $n$.

 Using the difference of squares, this simplifies to $(1+2+3+4+\cdots+n) \pi$. So, we are basically finding the smallest $n$ such that $\frac{n(n+1)}{2}>2023 \Rightarrow n(n+1) > 4046$. Since $60(61) > 60^2=3600$, the only option higher than $60$ is $\boxed{\textbf{(E) } 64}$.

 



\textbf{Solution 2}

After first observing the problem, we can work out a few of the areas.

 1st area = $4\pi-\pi = 3\pi$

 2nd area = $16\pi-9\pi = 7\pi$

 3rd area = $36\pi-25\pi = 11\pi$

 4th area = $64\pi-49\pi = 15\pi$

 We can see that the pattern is an arithmetic sequence with first term $3$ and common difference $4$. From here, we can start from the bottom of the answer choices and work our way up. As the question asks for the least number of circles needed total, we have to divide the number of total circles by 2.

 We can find the sum of the first $32$ terms of the arithmetic sequence by using the formula.

 The last term is: $3 + 4\cdot(32-1) = 127$. 

 Then, we can find the sum: $(3+127)/2 \cdot 32 = 65\cdot32 = 2080$. It is clear that $64$ works.

 The next answer choice is $60$, which we have to divide by 2 to get $30$.

 The last term is: $3 + 4\cdot(30-1) = 119$.

 The sum is: $(3+119)/2 \cdot 30 = 61\cdot30 = 1830$. This does not work.

 As answer choice $D$ does not work and $E$ does, we can conclude that the answer is $\boxed{\textbf{(E) } 64}$.

 



\textbf{Solution 3 (Similar to Solution 2)}

We can easily see that all of the answer choices are even, which helps us solve this problem a little.

 Let's just not consider the $\pi$, since it is not that important, so let's cancel that out.

 When we plug in 64, we get $64^2-63^2+62^2-61^2+\cdots +4^2-3^2+2^2-1^2$. By difference of squares, we get $1+2+3+\cdots+62+63+64$, which by the sum of an arithmetic sequence, is $\frac{64(64+1)}{2}$, which is $2080$.

 Similarly, we can use this for answer choice $D$, and we have $\frac{60(60+1)}{2}$ which is $1830$.

 So, we see that answer choice $D$ is too small to satisfy the requirements, so we conclude the answer is $\boxed{\textbf{(E) } 64}$.

 



\textbf{Solution 4}

We can consider making a table.

 If there is 1 circle, the area of the shaded region is 0$\pi$. (We can write this as 0$\pi$.)

 If there are 2 circles, the area of the shaded region is 3$\pi$.  (We can write this as (1+2)$\pi$).

 If there are 3 circles, the area of the shaded region is 3$\pi$.  (We can write this as (1+2)$\pi$).

 If there are 4 circles, the area of the shaded region is 10$\pi$.  (We can write this as (1+2+3+4)$\pi$).

 If there are 5 circles, the area of the shaded region is 10$\pi$.  (We can write this as (1+2+3+4)$\pi$).

 If there are 6 circles, the area of the shaded region is 21$\pi$.  (We can write this as (1+2+3+4+5+6)$\pi$).

 Now the pattern emerges. When $n$ is even, the area of the shaded region is $(1+2+3+\cdots+n)\pi$, or $\left( \frac{n(n+1)}{2} \right) \pi$. But remember that the problem stated that there are an even number of circles. So now we are solving the equation 
 \[\left( \frac{n(n+1)}{2} \right) \pi\ge2023\pi.\] Dividing by $\pi$ and multiplying by 2 on both sides, we get $n(n+1)\ge4046$. Now we can plug in the answer choices, and we start off with $60$ because it is the only answer choice that is a multiple of $10$. Plugging in we get $60\cdot61=3660$, and this is not quite yet more than $4046$. But only option $(\text{E})$ is bigger, so we know that the solution can only be $\boxed{\textbf{(E) } 64}$.

 (P.S. I did this in the test and I solved it in 2 minutes)

 



\textbf{Solution 5}

Denote $S$ the area of shaded region and $W$ the area of white region.

 $S > W$ and $S \approx W$ if $R$ is big.Therefore 
 \[\pi R^2 = S + W \approx 2S \ge 4046 \pi \implies R \approx 64.\]So we conclude the answer is $\boxed{\textbf{(E) } 64}$.

 \textbf{vladimir.shelomovskii@gmail.com, vvsss}

 



\textbf{Solution 6 (A Bit slow.)}

After first observing the problem, we can work out a few of the areas.

 1st area = $4\pi-\pi = 3\pi$

 2nd area = $16\pi-9\pi = 7\pi$

 3rd area = $36\pi-25\pi = 11\pi$

 4th area = $64\pi-49\pi = 15\pi$

 Notice that the area of the shaded region is $3\pi+7\pi+11\pi+15\pi...+(4n-1)\pi$ for any even number $n$.

 We can create a formula to compute this sum by doing:$3\pi+7\pi+11\pi+15\pi...+(4n-1)\pi$+$(4n-1)\pi+...15\pi+11\pi+7\pi+3\pi$to get$(4n+2)*\frac{n}{2} \to (2n^2+n)\pi$It is important to realize that each $n$ is $2$ circles, the unshaded and the shaded.So that means for each answer choice there are $2x$ as many circles. Then we just plug in values (strategically) to get that $2*(\frac{64}{2})^2+(\frac{64}{2})$ or $\boxed{\textbf{(E) } 64}$ as the answer.

 


	\item (\textit{AMC 10 10A 2023}) In a table tennis tournament every participant played every other participant exactly once. Although there were twice as many right-handed players as left-handed players, the number of games won by left-handed players was $40\%$ more than the number of games won by right-handed players. (There were no ties and no ambidextrous players.) What is the total number of games played?


 $\textbf{(A) }15\qquad\textbf{(B) }36\qquad\textbf{(C) }45\qquad\textbf{(D) }48\qquad\textbf{(E) }66$

 \textbf{Answer: }B



\textbf{Solution 1}

We know that the total amount of games must be the sum of games won by left and right handed players. Then, we can write $g = l + r$, and since $l = 1.4r$, $g = 2.4r$. Given that $r$ and $g$ are both integers, $g/2.4$ also must be an integer. From here we can see that $g$ must be divisible by 12, leaving only answers B and D. Now we know the formula for how many games are played in this tournament is $n(n-1)/2$, the sum of the first $n-1$ numbers. Now, setting 36 and 48 equal to the equation will show that two consecutive numbers must have a product of 72 or 96. Clearly, $72=8*9$, so the answer is $\boxed{\textbf{(B) }36}$.

 



\textbf{Solution 2}

First, we know that every player played every other player, so there's a total of $\dbinom{n}{2}$ games since each pair of players forms a bijection to a game. Therefore, that rules out D. Also, if the right-handed players won a total of $x$ games, the left-handed players must have won a total of $\dfrac{7}{5}x$ games, meaning that the total number of games played was $\dfrac{12}{5}x$. Thus, the total number of games must be divisible by $12$. Therefore leaving only answer choices B and D. Since answer choice D doesn't satisfy the first condition, the only answer that satisfies both conditions is $\boxed{\textbf{(B) }36}$

 


 A simpler way is to know that the result must be a number in the form $\dbinom{n}{2}$ games. This eliminates D. Then write a expression that equals the answer, $1.4x + x = n$ games total. Only $48$ and $36$ satisfy so $36$ is the answer which is $\boxed{\textbf{(B) }36}$




\textbf{Solution 3}

Let $r$ be the amount of games the right-handed won. Since the left-handed won $1.4r$ games, the total number of games played can be expressed as $(2.4)r$, or $12/5r$, meaning that the answer is divisible by 12. This brings us down to two answer choices, $B$ and $D$. We note that the answer is some number $x$ choose $2$. This means the answer is in the form $x(x-1)/2$. Since answer choice D gives $48 = x(x-1)/2$, and $96 = x(x-1)$ has no integer solutions, we know that $\boxed{\textbf{(B) }36}$ is the only possible choice.

 



\textbf{Solution 4 (Rigid solving, no cheese)}

Here is the rigid way to prove that $36$ is the only answer. Let the number of left-handed players be $n$, so the number of right-handed players is $2n$. The number of games won by the left-handed players comes in two ways: 

 
            \begin{itemize}
        \item The games played by two left-left pairs, which is $\frac{n(n-1)}{2}$, and
        \item The games played by left-right pairs, which we'll call $x$.
            \end{itemize}
            Note that $x\leq 2n^2,$ which is the total number of games played by left-right pairs. Using the same logic for right-right pairs and right-left pairs, we have that 
 \[\frac{\frac{n(n-1)}{2}+x}{\frac{2n(2n-1)}{2}+2n^2-x}=1.4,\] which gives 
 \[x=\frac{17n^2}{8}-\frac{3n}{8}.\] We know that $x\leq 2n^2$, applying that becomes
 \begin{align*} \frac{17n^2}{8}-\frac{3n}{8} &\le 2n^2 \\ 17n^2 - 3n &\le 16n^2 \\ n^2 - 3n &\le 0 \\ n &\le 3. \end{align*}
(We can safely divide by $n$ because it must be positive). So the total number of players $p$ can only be $3$, $6$, and $9$. 

 Since the total number of games $p(p-1)/2$  is $1.4+1 = \frac{12}{5}$ times a non-negative integer number of games won by righties, $p(p-1)$ must be a multiple of $12(2) = 24$. Among $\{3,6,9\}$, only $p = 9$ satisfies this condition,  so the total number of games is $\frac{(9)(8)}{2} = \boxed{\textbf{(B)}\ 36}.$




\textbf{Solution 5 (Cheese)}

If there are $n$ players, the total number of games played must be $\binom{n}{2}$, so it has to be a triangular number. The ratio of games won by left-handed to right-handed players is $7:5$, so the number of games played must also be divisible by $12$. Finally, we notice that only $\boxed{\textbf{(B) }36}$ satisfies both of these conditions.

 



\textbf{Solution 5.1 (Slightly Slower)}

We know that the number of games played must be a triangular number, which elimiates ${\textbf{(D)}48}$, then we see that there are twice as many right-handed players as left-handed ones. Which means that the total number of players must be divisible by $3$, and since each player can't play themselves, we can add $1$ to the largest number in the sequence, since $1 + 2 + 3 + \dots + 8$ is the only one which the amount of players is divisble by $3(9 players)$, $\boxed{\textbf{(B) }36}$ is the answer

 



\textbf{Solution 6 (Cheese 2)}

Call the number of games x, the number of games won by right handed player R, and by left handed players L. L+R=x and L=1.4 R. Therefore, x/2.4 must be an integer, which leaves only 36 and 48. Then, note x must be half the product of two consecutive numbers (the number of possible combinations of two people) and we can eliminate 48.

 
How is this cheese... it is literally the first solution but scaled down

 -Shreyansh

 EXACTLY... the first solution is also cheese.

 -AndrewZhong2012

 cheddar

 -PerseverePlayer

 



\textbf{Solution 7}

Let the number of people be $n$. This means that there will be $\frac{n(n-1)}{2}$ games played. We can now eliminate $\textbf{(D)}$ because it's not a triangular number. The problem states that there are twice as many right handed people as left handed people with no ambidextrous people. Based on this, we can say that the number of people (or $n$) has to be divisible by 3. Because the number of games won by left-handed people is 40\% more than the number of games won by right-handed people, we can set the ratio between the number of games won by left-handed to right-handed as $\frac{7}{5}$. The sum of the numerator and the denominator is 12. This means that the number of games must be divisible by 12. Thus, the only answer choice that fits the conditions above is answer $\boxed{\textbf{(B) }36}$.

 



\textbf{Solution 8(Pairing) - easiest noncheese}

Case 1: in-group pairing

 
In every pair where both players use the same hand, the win rate is always \(50\%\). This is because the total number of such pairs is even — any factorial other than the trivial case \(1!\) is even, and \(1\) is insufficient to form a pair. Thus, no player is left out, and in each pair exactly one player wins and one loses, giving a win‑loss ratio of \(1:1\). In other words, this case completely cancels out due to its balanced win‑loss ratio.

 
Case 2: cross-group pairing

 
For convenience, let left‑handed players be denoted by \(L\) and right‑handed players by \(R\). The problem gives the ratio \(R : L = 1 : 2\). This means each \(L\) plays \(2L\) times as many games as an \(R\). From this, the total number of cross‑group games is \(2L \times L = 2L^2\).

 The problem also states that the ratio of \(L\) wins to \(R\) wins is \(7:5\), so the total number of wins in this case must be divisible by \(5 + 7 = 12\). We therefore require \(2L^2\) to be divisible by \(12\). Simplifying gives \(L^2 \mid 6\). The smallest perfect square divisible by \(6\) is \(\boxed{\textbf{(B)}\ 36}\).

 



\textbf{Solution 9(Subgroups)}

Let L and R be the number of left-handed and right-handed players, respectively. The total number of games played in the tournament can be broken down into three distinct parts based on the matchups: games between two left-handed players, games between two right-handed players, and games between a left-handed and a right-handed player. This gives the total game count formula:

 
 \[\text{Total Games} = LR + \frac{L(L-1)}{2} + \frac{R(R-1)}{2}\]

 Every game played results in exactly one win. In the matches where left-handed players play each other, a left-handed player must win. This guarantees $\frac{L(L-1)}{2}$ wins for the left-handed group. Similarly, matches where right-handed players play each other guarantee $\frac{R(R-1)}{2}$ wins for the right-handed group. 

 For the LR games played between a left-handed player and a right-handed player, let W be the number of games won by the left-handed players. The remaining (LR - W) games in this section must be won by the right-handed players. We can now write expressions for the total wins of each group:

 
 \[\text{Total Left Wins } (W_L) = \frac{L(L-1)}{2} + W\]
 \[\text{Total Right Wins } (W_R) = \frac{R(R-1)}{2} + LR - W\]

 The problem states that there are twice as many right-handed players as left-handed players, so we substitute R = 2L into our equations:

 
 \[W_L = \frac{L(L-1)}{2} + W\]
 \[W_R = \frac{2L(2L-1)}{2} + L(2L) - W = L(2L-1) + 2L^2 - W = 4L^2 - L - W\]

 We are given that left-handed players won 40\% more games than right-handed players, which means $W_L = 1.4 \times W_R$. Setting up the equation:

 
 \[\frac{L(L-1)}{2} + W = 1.4(4L^2 - L - W)\]

 To clear the fraction and decimal, we multiply the entire equation by 10:

 
 \[5L(L-1) + 10W = 14(4L^2 - L - W)\]
 \[5L^2 - 5L + 10W = 56L^2 - 14L - 14W\]

 Grouping the W terms on the left side and the L terms on the right side yields:

 
 \[24W = 51L^2 - 9L\]

 Dividing the entire equation by 3 simplifies it to:

 
 \[8W = 17L^2 - 3L \implies W = \frac{17L^2 - 3L}{8}\]

 Since W represents a physical number of wins, it must be a non-negative integer. Additionally, W cannot be greater than the total number of inter-group games, meaning $W \le 2L^2$. Testing small integer values for L:

 
            \begin{itemize}
        \item If L=1: $W = \frac{14}{8}$, which is not an integer.
        \item If L=2: $W = \frac{62}{8}$, which is not an integer.
        \item If L=3: $W = \frac{153 - 9}{8} = \frac{144}{8} = 18$. This is a valid integer, and it satisfies $18 \le 2(3)^2 = 18$.
            \end{itemize}
            With L=3, the number of right-handed players is R = 6, giving a total of 9 players. Substituting these into the original total games formula:

 
 \[\text{Total Games} = (3)(6) + \frac{3(2)}{2} + \frac{6(5)}{2} = 18 + 3 + 15 = \boxed{\textbf{(B)}\ 36}\]



	\item (\textit{AMC 10 10A 2023}) Let $ABCD$ be a rectangle with $AB = 30$ and $BC = 28$. Point $P$ and $Q$ lie on $\overline{BC}$ and $\overline{CD}$ respectively so that all sides of $\triangle{ABP}, \triangle{PCQ},$ and $\triangle{QDA}$ have integer lengths. What is the perimeter of $\triangle{APQ}$?


 $\textbf{(A) } 84 \qquad \textbf{(B) } 86 \qquad \textbf{(C) } 88 \qquad \textbf{(D) } 90 \qquad \textbf{(E) } 92$

 \textbf{Answer: }A



\textbf{Solution 1 (Testing Pythagorean Triples)}


                    \begin{figure}[H]
                        \centering
                        \includegraphics[width=0.5\textwidth]{Solution/AMC_10_10A_2023_17_Solution_Figure_1.png}
                    \end{figure}
                    

 We know that all side lengths are integers, so we can test Pythagorean triples for all triangles. 

 First, we focus on $\triangle{ABP}$. The length of $AB$ is $30$, and the possible (small enough) Pythagorean triples $\triangle{ABP}$ can be are $(3, 4, 5), (5, 12, 13), (8, 15, 17),$ where the length of the longer leg is a factor of $30$. Testing these, we get that only $(8, 15, 17)$ is a valid solution. Thus, we know that $BP = 16$ and $AP = 34$. 

 Next, we move on to $\triangle{QDA}$. The length of $AD$ is $28$, and the small enough triples are $(3, 4, 5)$ and $(7, 24, 25)$. Testing again, we get that $(3, 4, 5)$ is our triple. We get the value of $DQ = 21$, and $AQ = 35$.

 We know that $CQ = CD - DQ$ which is $9$, and $CP = BC - BP$ which is $12$. $\triangle{CPQ}$ is therefore a right triangle with side length ratios ${3, 4, 5}$, and the hypotenuse is equal to $15$.$\triangle{APQ}$ has side lengths $34, 35,$ and $15,$ so the perimeter is equal to $34 + 35 + 15 = \boxed{\textbf{(A) } 84}.$




\textbf{Solution 2}

Let $BP=y$ and $AP=z$. We get $30^{2}+y^{2}=z^{2}$. Subtracting $y^{2}$ on both sides, we get $30^{2}=z^{2}-y^{2}$. Factoring, we get $30^{2}=(z-y)(z+y)$. Since $y$ and $z$ are integers, both $z-y$ and $z+y$ have to be even or both have to be odd. We also have $y<31$. We can pretty easily see now that $z-y=18$ and $z+y=50$. Thus, $y=16$ and $z=34$. We now get $CP=12$. We do the same trick again. Let $DQ=a$ and $AQ=b$. Thus, $28^{2}=(b+a)(b-a)$. We can get $b+a=56$ and $b-a=14$. Thus, $b=35$ and $a=21$. We get $CQ=9$ and by the Pythagorean Theorem, we have $PQ=15$. We get $AP+PQ+AQ=34+15+35=84$. Our answer is $\boxed{\textbf{(A) } 84}.$

 If you want to see a video solution on this solution, look at Video Solution 1.

 -paixiao (minor edits-HW73)

 


	\item (\textit{AMC 10 10A 2023}) A rhombic dodecahedron is a solid with $12$ congruent rhombus faces. At every vertex, $3$ or $4$ edges meet, depending on the vertex. How many vertices have exactly $3$ edges meet?


 $\textbf{(A) }5\qquad\textbf{(B) }6\qquad\textbf{(C) }7\qquad\textbf{(D) }8\qquad\textbf{(E) }9$

 \textbf{Answer: }D



\textbf{Solution 1}

Note Euler's formula where $\text{Vertices}+\text{Faces}-\text{Edges}=2$. There are $12$ faces. There are $24$ edges, because there are 12 faces each with four edges and each edge is shared by two faces. Now we know that there are $2-12+24=14$ vertices. Now note that the sum of the degrees of all the points is $48$(the number of edges times 2). Let $x=$ the number of vertices with $3$ edges. Now we know $\frac{3x+4(14-x)}{2}=24$. Solving this equation gives $x = \boxed{\textbf{(D) }8}$. ~aiden22gao ~zgahzlkw (LaTeX) ~ESAOPS (Simplified)

 



\textbf{Solution 2}

Let $x$ be the number of vertices with 3 edges, and $y$ be the number of vertices with 4 edges. Since there are $\frac{4*12}{2}=24$ edges on the polyhedron, we can see that $\frac{3x+4y}{2}=24$. Then, $3x+4y=48$. Notice that by testing the answer choices, (D) is the only one that yields an integer solution for $y$. Thus, the answer is $\boxed{\textbf{(D) }8}$.

 



\textbf{Solution 3}

With $12$ rhombi, there are $4\cdot12=48$ total boundaries. Each edge is used as a boundary twice, once for each face on either side. Thus we have $\dfrac{48}2=24$ total edges.

 Let $A$ be the number of vertices with $3$ edges (this is what the problem asks for) and $B$ be the number of vertices with $4$ edges. We have $3A + 4B = 48$.

 Euler's formula states that, for all convex polyhedra, $V-E+F=2$. In our case, $V-24+12=2\implies V=14.$ We know that $A+B$ is the total number of vertices as we are given that all vertices are connected to either $3$ or $4$ edges. Therefore, $A+B=14.$

 We now have a system of two equations. Solving the system yields $A=\boxed{\textbf{(D) }8}$.

 Even without Euler's formula, we can do a bit of answer guessing. From $3A+4B=48$, we take mod $4$ on both sides.

 
 \[3A+4B\equiv48\pmod{4}\]
 \[3A\equiv0\pmod{4}\]

 We know that $3A$ must be divisible by $4$. We know that the factor of $3$ will not affect the divisibility by $4$ of $3A$, so we remove the $3$. We know that $A$ is divisible by $4$. Checking answer choices, the only one divisible by $4$ is indeed $A=\boxed{\textbf{(D) }8}$.

 



\textbf{Solution 4}

Note that Euler's formula is $V+F-E=2$. We know $F=12$ from the question. We also know $E = \frac{12 \cdot 4}{2} = 24$ because every face has $4$ edges and every edge is shared by $2$ faces. We can solve for the vertices based on this information.

 Using the formula we can find:
 \[V + 12 - 24 = 2\]
 \[V = 14\]Let $t$ be the number of vertices with $3$ edges and $f$ be the number of vertices with $4$ edges. We know $t+f = 14$ from the question and $3t + 4f = 48$. The second equation is because the total number of points is $48$ because there are 12 rhombuses of $4$ vertices. Now, we just have to solve a system of equations.
 \[3t + 4f = 48\]
 \[3t + 3f = 42\]
 \[f = 6\]
 \[t = 8\]Our answer is simply just $t$, which is $\boxed{\textbf{(D) }8}$ ~musicalpenguin

 



\textbf{Solution 5}

Each of the twelve rhombi has two pairs of angles across from each other that must be congruent. If both pairs of angles occur at $4$-point intersections, we have a grid of squares. If both occur at $3$-point intersections, we would have a cube with six square faces. Therefore, two of the points must occur at a $3$-point intersection and two at a $4$-point intersection.

 Since each $3$-point intersection has $3$ adjacent rhombuses, we know the number of $3$-point intersections must equal the number of $3$-point intersections per rhombus times the number of rhombuses over $3$. Since there are $12$ rhombuses and two $3$-point intersections per rhombus, this works out to be:

 $\frac{2\cdot12}{3}$

 Hence: $\boxed{\textbf{(D) }8}$ ~hollph27 ~Minor edits by FutureSphinx

 



\textbf{Solution 6 (Based on previous knowledge)}

Note that a rhombic dodecahedron is formed when a cube is turned inside out (as seen here), thus there are 6 4-vertices (corresponding to each face of the cube) and 8 3-vertices (corresponding to each corner of the cube). Thus the answer is $\boxed{\textbf{(D) }8}$




\textbf{Solution 7 (Dual)}

Note that a rhombic dodecahedron is the dual of a cuboctahedron. A cuboctahedron has $8$ triangular faces, which correspond to $\boxed{\textbf{(D) }8}$ vertices on a rhombic dodecahedron that have $3$ edges.

 



\textbf{Solution 8 (Uses Solution 2)}

Continue through solution 2 until you get $3x+4y=48$. Then we have that $3x = 48-4y = 4(12-y) = \text{multiple of 4}$ so $x$ is a multiple of $4.$ The only multiple of 4 answer choice is $8 = \boxed{(D)}$



	\item (\textit{AMC 10 10A 2023}) The line segment formed by $A(1, 2)$ and $B(3, 3)$ is rotated to the line segment formed by $A'(3, 1)$ and $B'(4, 3)$ about the point $P(r, s)$. What is $|r-s|$?


 $\textbf{(A) } \frac{1}{4} \qquad \textbf{(B) } \frac{1}{2} \qquad \textbf{(C) } \frac{3}{4}   \qquad \textbf{(D) } \frac{2}{3} \qquad   \textbf{(E) } 1$

 \textbf{Answer: }E



\textbf{Solution 1}

Due to rotations preserving an equal distance, we can bash the answer with the distance formula. $D(A, P) = D(A', P)$, and $D(B, P) = D(B',P)$.Thus we will square our equations to yield: $(1-r)^2+(2-s)^2=(3-r)^2+(1-s)^2$, and $(3-r)^2+(3-s)^2=(4-r)^2+(3-s)^2$.Canceling $(3-s)^2$ from the second equation makes it clear that $r$ equals $3.5$.

 Substituting will yield 

 
 \begin{align*}(2.5)^2+(2-s)^2 &= (-0.5)^2+(1-s)^2 \\ 6.25+4-4s+s^2 &= 0.25+1-2s+s^2 \\ 2s &= 9 \\ s &=4.5 \\ \end{align*}
.

 Now $|r-s| = |3.5-4.5| = \boxed{\textbf{(E) } 1}$.

 -Antifreeze5420

 



\textbf{Solution 2}

Due to rotations preserving distance, we have that $BP = B^\prime P$, as well as $AP = A^\prime P$. From here, we can see that P must be on the perpendicular bisector of $\overline{BB^\prime}$ due to the property of perpendicular bisectors keeping the distance to two points constant.

 From here, we proceed to find the perpendicular bisector of $\overline{BB^\prime}$. We can see that this is just a horizontal line segment with midpoint at $(3.5, 3)$. This means that the equation of the perpendicular bisector is $x = 3.5$.

 Similarly, we find the perpendicular bisector of $\overline{AA^\prime}$. We find the slope to be $\frac{1-2}{3-1} = -\frac12$, so our new slope will be $2$. The midpoint of $A$ and $A^\prime$ is $(2, \frac32)$, which we can use with our slope to get another equation of $y = 2x - \frac52$.

 Now, point P has to lie on both of these perpendicular bisectors, meaning that it has to satisfy both equations. Plugging in the value of $x$ we found earlier, we find that $y=4.5$. This means that $|r - s| = |3.5 - 4.5| = \boxed{\textbf{(E) } 1}$.

 -DEVSAXENA

 



\textbf{Solution 3 (Coordinate Geometry)}

To find the center of rotation, we find the intersection point of the perpendicular bisectors of $\overline{AA^\prime}$ and $\overline{BB^\prime}$. (this is because $\overline{PA}$ is equal to $\overline{PA^\prime}$ and same with $B$ and $B^\prime$

 We can find that the equation of the line $\overline{AA^\prime}$ is $y = -\frac{1}{2}x + \frac{5}{2}$, and that the equation of the line $\overline{BB^\prime}$ is $y = 3$. 

 When we solve for the perpendicular bisector of  $y = -\frac{1}{2}x + \frac{5}{2}$, we determine that it has a slope of 2, and it runs through $(2, 1.5)$. Plugging in $1.5 = 2(2)-n$, we get than $n = \frac{5}{2}$. Therefore our perpendicular bisector is $y=2x-\frac{5}{2}$. Next, we solve for the perpendicular of $y = 3$. We know that it has an undefined slope, and it runs through $(3.5, 3)$. We can determine that our second perpendicular bisector is $x = 3.5$.

 Setting the equations equal to each other, we get $y = \frac{9}{2}$. Therefore, $|r - s| = |3.5 - 4.5| = \boxed{\textbf{(E) } 1}$.

 
                    \begin{figure}[H]
                        \centering
                        \includegraphics[width=0.5\textwidth]{Solution/AMC_10_10A_2023_19_Solution_Figure_1.png}
                    \end{figure}
                    




\textbf{Solution 4}

We use the complex numbers approach to solve this problem.Denote by $\theta$ the angle that $AB$ rotates about $P$ in the counterclockwise direction.

 Thus, $A' - P = e^{i \theta} \left( A - P \right)$ and $B' - P = e^{i \theta} \left( B - P \right)$.

 Taking ratio of these two equations, we get
 \[ \frac{A' - P}{A - P} = \frac{B' - P}{B - P} . \]

 By solving this equation, we get $P = \frac{7}{2} + i \frac{9}{2}$.Therefore, $|s-t| = \left| \frac{7}{2} - \frac{9}{2} \right| = \boxed{\textbf{(E) 1}}$.

 



\textbf{Solution 5 (Overkill Trigonometry)}

First, we need to find the coordinates of point \( P(r, s) \) such that rotating the line segment from \( A(1,2) \) to \( B(3,3) \) about \( P \) maps it to the line segment from \( A'(3,1) \) to \( B'(4,3) \). After finding \( r \) and \( s \), we will compute \( |r - s| \).

 When a point \( (x, y) \) is rotated about another point \( (r, s) \) by an angle \( \theta \), the new coordinates \( (x', y') \) are given by:

 
\begin{cases}x' = r + (x - r) \cos \theta - (y - s) \sin \theta \\y' = s + (x - r) \sin \theta + (y - s) \cos \theta\end{cases}

 
We apply the rotation to points \( A \) and \( B \):

 For \( A(1,2) \) mapping to \( A'(3,1) \):

 
\begin{cases}3 = r + (1 - r) \cos \theta - (2 - s) \sin \theta \quad \text{(1)} \\1 = s + (1 - r) \sin \theta + (2 - s) \cos \theta \quad \text{(2)}\end{cases}

 
For \( B(3,3) \) mapping to \( B'(4,3) \):

 
\begin{cases}4 = r + (3 - r) \cos \theta - (3 - s) \sin \theta \quad \text{(3)} \\3 = s + (3 - r) \sin \theta + (3 - s) \cos \theta \quad \text{(4)}\end{cases}

 
Now, we subtract equations to eliminate variables:

 Subtract equation (1) from equation (3):

 
 \[(4 - 3) = [r + (3 - r) \cos \theta - (3 - s) \sin \theta] - [r + (1 - r) \cos \theta - (2 - s) \sin \theta]\]

 Simplifying yields:

 
 \[1 = 2 \cos \theta - \sin \theta \quad \text{(A)}\]

 Similarly, subtract equation (2) from equation (4):

 
 \[(3 - 1) = [s + (3 - r) \sin \theta + (3 - s) \cos \theta] - [s + (1 - r) \sin \theta + (2 - s) \cos \theta]\]

 Simplifying further:

 
 \[2 = 2 \sin \theta + \cos \theta \quad \text{(B)}\]

 Now, using equations (A) and (B):

 
\begin{cases}2 \cos \theta - \sin \theta = 1 \\\cos \theta + 2 \sin \theta = 2\end{cases}

 
We now express \( \cos \theta \) from the second equation:

 
 \[\cos \theta = 2 - 2 \sin \theta\]

 Substituting this into the first equation, we get:

 
 \[2(2 - 2 \sin \theta) - \sin \theta = 1 \\ 4 - 4 \sin \theta - \sin \theta = 1 \\ -5 \sin \theta = -3 \\ \sin \theta = \frac{3}{5}\]

 Now, we find $\cos \theta$ as follows:

 
 \[\cos \theta = 2 - 2 \left( \frac{3}{5} \right) = \frac{4}{5}\]

 Substituting $\cos \theta$ and $\sin \theta$ back into equations (1) and (2):

 Equation (1):

 
 \[3 - r = (1 - r) \left( \frac{4}{5} \right) - (2 - s) \left( \frac{3}{5} \right)\]

 Multiply both sides by 5:

 
 \[5(3 - r) = 4(1 - r) - 3(2 - s)\]

 Simplify:

 
 \[15 - 5r = 4 - 4r - 6 + 3s \\ 15 - 5r = -2 - 4r + 3s \\ -5r + 4r - 3s = -2 - 15 \\ -r - 3s = -17 \quad \text{(C)}\]

 Equation (2):

 
 \[1 - s = (1 - r) \left( \frac{3}{5} \right) + (2 - s) \left( \frac{4}{5} \right)\]

 Multiply both sides by 5:

 
 \[5(1 - s) = 3(1 - r) + 4(2 - s)\]

 Simplifying more:

 
 \[5 - 5s = 3 - 3r + 8 - 4s \\ 5 - 5s = 11 - 3r - 4s \\ -5s + 4s + 3r = 11 - 5 \\ 3r - s = 6 \quad \text{(D)}\]

 From equations (C) and (D):

 
 \[\begin{cases} -r - 3s = -17 \\ 3r - s = 6 \end{cases}\]

 We then multiply the first equation by 3:

 
 \[-3r - 9s = -51\]

 Add this to the second equation multiplied by 9:

 
 \[27r - 9s = 54\]

 Subtract:

 
 \[(27r - 9s) - (-3r - 9s) = 54 - (-51) \\ 30r = 105 \\ r = \frac{105}{30} = \frac{7}{2} = 3.5\]

 Now find \( s \):

 
 \[3(3.5) - s = 6 \\ 10.5 - s = 6 \\ s = 4.5\]

 Therefore, we have the following:

 
 \[|r - s| = |3.5 - 4.5| = |-1| = \boxed{\textbf{(E) } 1}.\]



	\item (\textit{AMC 10 10A 2023}) Each square in a $3\times3$ grid of squares is colored red, white, blue, or green so that every $2\times2$ square contains one square of each color. One such coloring is shown on the right below. How many different colorings are possible? 


 
                    \begin{figure}[H]
                        \centering
                        \includegraphics[width=0.5\textwidth]{Problem/AMC_10_10A_2023_20_Problem_Figure_1.png}
                    \end{figure}
                    
$\textbf{(A) }24\qquad\textbf{(B) }48\qquad\textbf{(C) }60\qquad\textbf{(D) }72\qquad\textbf{(E) }96$

 \textbf{Answer: }D



\textbf{Solution 1}

Let a "tile" denote a \(1 \times 1\) square, and a "square" refer to a \(2 \times 2\) square.

 We have \(4! = 24\) possible ways to fill out the top-left square. Next, we fill out the bottom-right corner tile. In the bottom-right square, one corner is already filled (the central tile) from our initial coloring), so we have 3 color options remaining for this. 

 Now considering the remaining tiles, all of these only have one way to be filled (Try it yourself if you don’t believe).

 For example, the right tile of the middle row is part of two squares: the top-right and the bottom-right. Among these squares, 3 colors have already been used, leaving us with only 1 remaining option. Similarly, every other remaining square has only one available option for coloring. 

 Thus, the total number of ways is \(3 \times 4! = \boxed{\textbf{(D) }72}\).

 A quick version of this method is used when you use for example B as the top left corner block. You can see that $B$ in the top left corner has $6$ possible ways. Now, see that there are $3$ possible corners for a center cell and there are $4$ possible center cells. We get $6$ times $4$ times $3 = D$, the answer is \( \boxed{\textbf{(D) }72}\). (I'm not very good at using latex)~breakingbread~latex done by pungent\_muskrat

 



\textbf{Solution 2 (Circular Counting)}

Recall the circular gap theorem
 \[C_n(A) = \frac{n}{k} \binom{n-k-1}{k-1}.\]We have nine slots and a maximum of 4 colors cannot be next to one another (if you aren't convinced, try plugging in any number greater than 4 and see what happens to the formula; it becomes invalid, and is therefore 0). Therefore, we take the cases $n=9$ and $k \in \{1,2,3,4\}$.

 Case 1: $k \in \{1\}$

 This is simply just $\frac{9}{1} \binom{9-1-1}{1-1} = 9$ ways.

 Case 2: $k \in \{2\}$

 This is just $\frac{9}{2} \binom{9-2-1}{2-1} = \frac{9}{2} \cdot 6 = 27$ ways.

 Case 3: $k \in \{3\}$

 This is $\frac{9}{3} \binom{9-3-1}{3-1} = 3 \cdot 10 = 30$ ways.

 Case 4: $k \in \{4\}$

 This is $\frac{9}{4} \binom{9-4-1}{4-1} = \frac{9}{4} \cdot 4 = 9$ ways.

 The number of ways to arrange the colors such that no same color is next to each other at least once is $9+27+30+9 = 75$.

 However, it is clear that every color (excluding the center) is only used a maximum of three times, so we apply the formula for $k \in \{1,2,3\}$ and $n = 3$ (three total slots it can go to). Notice that $k \in \{2,3\}$ is just 0, so therefore we just consider $k \in \{1\}$ and get
 \[\frac{3}{1} \binom{3-1-1}{1-1} = 3\]

 Our final answer is just $75 - 3 =$ $\boxed{\textbf{(D) }72}$




\textbf{Solution 3}


                    \begin{figure}[H]
                        \centering
                        \includegraphics[width=0.5\textwidth]{Solution/AMC_10_10A_2023_20_Solution_Figure_1.png}
                    \end{figure}
                    We can split this problem into $2$ cases as shown in solution 1. We can swap a set of equal colors for another set of equal colors to create a new square.

 Square 1:The first square can be rotated to create another square so we have to multiply the number of arrangements by $2$. We have $4! = 24$ arrangements without rotating and $24\cdot 2 = 48$ arrangements in total for the first square.

 Square 2:There are $4! = 24$ ways to arrange the colors.

 In total, we have $48 + 24 = \boxed{\textbf{(D) }72}$ arrangements.

 



\textbf{Solution 4}

Let’s call the top-right corner color A, the top-middle color B, the top-right color C, and so on, with color D being the middle row, and right corner square, and color G being the bottom-left square’s color.

 WLOG A=Red, B=White, D=Blue, and E=Green. We will now consider squares C and F’s colors.Case 1 : C=Red and F=Blue

 In this case, we get that G and H have to be Red and White in some order, and the same for H and I. We can color this in 2 ways.Case 2 : C=Blue and F=Red

 In this case, one of G and H needs to be White and Red, and  H and I needs to be White and Blue. There is 1 way to color this.In total, we get $24*(2+1)=72$ ways to color the grid. $\boxed{\textbf{(D) }72}$.

 -paixiao

 -jkim0656 (minor formatting and latex)

 



\textbf{Solution 5}

We will choose colors step-by-step:

 1. There are $4$ ways to choose a color in the center.

 2. Then we select any corner and there would be $3$ ways to choose a color as we can't use the same color as the one in the center.

 3. Consider the $2\times 2$ square that contains the center and the corner we have selected. For the other $2$ squares, there are $2$ ways to choose colors.

 4. Now, consider how many configurations it makes sense to construct the $2\times 2$ square opposite to the corner we have selected using the $2$ other $2\times 2$ squares, and we get $3$ configurations.

 Finally, the answer is $4 \cdot 3 \cdot 2 \cdot 3 = \boxed{\textbf{(D) }72}$




\textbf{Solution 6 (possibly bad)}


                    \begin{figure}[H]
                        \centering
                        \includegraphics[width=0.5\textwidth]{Solution/AMC_10_10A_2023_20_Solution_Figure_2.png}
                    \end{figure}
                    

 
Note that we could fill the 3 by 3 square with numbers of choices, rather than letters or color names. The top-left corner receives a 3 because there are 3 total choices to choose from: R, G and B. The squares bordering them has values of 2 and 1, regardless of order. 2 indicates that the small square can have any color excluding the one existing color, 1 indicates the remaining color of the 2 by 2 square. Finally, the middle square receives 3, since the first 2 by 2 square has RGB already, and it does not matter what color it has. Moving onto the next 2 by 2 square, we see that there are already 2 decided colors, so the other squares must be 2 and 1, again, regardless of the order. The same applys to the third 2 by 2 square, and finally the last square has the value of one, as it is the only square left. Multiplying the numbers, $2\times2\times2\times3\time3\times3$ = $\boxed{\textbf{(D) }72}$ (Similar to Solution 1)

 -MEZE\_RUNwhy are there only 3 choices for the top right corner? you havent placed any colors yet so there should be 4 choices which would make your answer incorrect

 



\textbf{Solution 7}


                    \begin{figure}[H]
                        \centering
                        \includegraphics[width=0.5\textwidth]{Solution/AMC_10_10A_2023_20_Solution_Figure_3.png}
                    \end{figure}
                    
\begin{verbatim}
Let's name the cells $A,B,C,D,E,F,G,H,I$ from the top left to the bottom right. 
\end{verbatim}

Case 1.  Cell $B$ and cell $H$ have the same color. The middle one cell $E$ has $4$ choices, cell $B$ has $3$ choices, then cell $E$ has $2$ choices and cell $F$ has $2$ choices, this gives $4\cdot 3\cdot 2\cdot 2=48$ ways.

 Case 2. Cell $B$ and cell $H$ have different colors. The middle one cell $E$ has $4$ choices, cell $B$ has $3$ choices, cell $H$ has $2$ choices, then cell $D$ and $F$ each can only have one choice (different from $B,E,H).$ This gives$4\cdot 3\cdot 2=24$ ways. 

 The answer is thus $48+24=72,$ corresponding to answer choice $\boxed{D}$




\textbf{Solution 8}

Consider Cases based on the color of the top left and bottom left square. 

 Case 1: The top left and bottom left squares are the same:

 We have $4$ ways to assign the top and bottom squares, lets say it is red. Once that is done, we can assign the center square in 3 ways, since it can be anything except the color of the top and bottom squares (Lets say it is white). We can then assign the bottom middle square any color except the center and bottom left square, so there is two ways to do that (Lets say its blue). From there, you can assign the other squares in the first two columns in only one way (Top middle = Blue, Middle Left = Green). Then, the right column can go in one of two ways: Green, Red, Green, or Red, Green, Red, for a total of two ways. Taking this altogether gives $4 * 3 * 2 * 2 = 48$

 Case 2: The top left and bottom left squares are different:

 We have $4 * 3 = 12$ ways to assign the top and bottom squares, let them be Red and Green respectively. Then the square between them can be assigned in two ways. We realize how the  middle square can only be filled in one way (White, because it can't be Blue, Green, or Red). From there, the entire second column can only be filled in one way (From top to bottom, it is Green, White, Red). The same story goes for the last column, it can only be filled in one way (From top to bottom it is Red, Blue, Green). That gives $4 * 3 * 2 = 24$ ways. 

 Thus, the answer is  $48+24=72,$ which is answer choice $\boxed{\textbf{(D) }72}$~latex work done by pungent\_muskrat$

 



\textbf{Solution 9}

We can first choose the center square. There are $4$ possibilities: Red, Blue, Green, or White. WLOG, lets assume that the center is white. Next, there are $3$ possibilities for the top middle square: Red, Blue, or Green. WLOG, lets assume that it is Red. We now have two cases:

 CASE $1$: The first and last column are in the form BG... and BG... or GB... and GB...:We can see that there are two possibilities for this case: the bottom row is RBR or BRB. The total number of ways for this case is $4\cdot3\cdot2\cdot2$.CASE $2$: The first and last column are in the form BG... and GB... or GB... and BG...We can now see that the bottom middle cell HAS to be a red (if it isn't, either the bottom left or bottom right 2x2 squares will have  duplicate colors). After we choose the bottom middle, the rest of the cells are automatically chosen. Thus, the total number of ways for this case is $4\cdot3\cdot2\cdot1$.Thus, our final answer is the sum of case 1 and case 2, which is $48 + 24 = 72$, or $\boxed{\textbf{(D) }72}$.~ Irfans123

 


	\item (\textit{AMC 10 10A 2023}) Let $P(x)$ be the unique polynomial of minimal degree with the following properties:


 
            \begin{itemize}
        \item $P(x)$ has a leading coefficient $1$,
            \end{itemize}
            
            \begin{itemize}
        \item $1$ is a root of $P(x)-1$,
            \end{itemize}
            
            \begin{itemize}
        \item $2$ is a root of $P(x-2)$,
            \end{itemize}
            
            \begin{itemize}
        \item $3$ is a root of $P(3x)$, and
            \end{itemize}
            
            \begin{itemize}
        \item $4$ is a root of $4P(x)$.
            \end{itemize}
            The roots of $P(x)$ are integers, with one exception. The root that is not an integer can be written as $\frac{m}{n}$, where $m$ and $n$ are relatively prime integers. What is $m+n$?


 $\textbf{(A) }41\qquad\textbf{(B) }43\qquad\textbf{(C) }45\qquad\textbf{(D) }47\qquad\textbf{(E) }49$

 \textbf{Answer: }D



\textbf{Solution 1}

From the problem statement, we find $P(2-2)=0$, $P(9)=0$ and $4P(4)=0$. Therefore, we know that $0$, $9$, and $4$ are roots. So, we can factor $P(x)$ as $x(x - 9)(x - 4)(x - a)$, where $a$ is the unknown root. Since $P(x) - 1 = 0$, we plug in $x = 1$ which gives $1(-8)(-3)(1 - a) = 1$, so $24(1 - a) = 1 \implies 1 - a = 1/24 \implies a = 23/24$. Therefore, our answer is $23 + 24 =\boxed{\textbf{(D) }47}$




\textbf{Solution 2}

We proceed similarly to solution one. We get that $x(x-9)(x-4)(x-a)=1$. Expanding, we get that $x(x-9)(x-4)(x-a)=x^4-(a+13)x^3+(13a+36)x^2-36ax$. We know that $P(1)=1$, so the sum of the coefficients of the cubic expression is equal to one. Thus $1-(a+13)+(13a+36)-36a=1$. Solving for $a$, we get that $a = 23/24$. Therefore, our answer is $23 + 24 =\boxed{\textbf{(D) }47}$



	\item (\textit{AMC 10 10A 2023}) Circle $C_1$ and $C_2$ each have radius $1$, and the distance between their centers is $\frac{1}{2}$. Circle $C_3$ is the largest circle internally tangent to both $C_1$ and $C_2$. Circle $C_4$ is internally tangent to both $C_1$ and $C_2$ and externally tangent to $C_3$. What is the radius of $C_4$? 


 
                    \begin{figure}[H]
                        \centering
                        \includegraphics[width=0.5\textwidth]{Problem/AMC_10_10A_2023_22_Problem_Figure_1.png}
                    \end{figure}
                    
$\textbf{(A) } \frac{1}{14} \qquad \textbf{(B) } \frac{1}{12} \qquad \textbf{(C) } \frac{1}{10} \qquad \textbf{(D) } \frac{3}{28} \qquad \textbf{(E) } \frac{1}{9}$

 \textbf{Answer: }D



\textbf{Solution}


                    \begin{figure}[H]
                        \centering
                        \includegraphics[width=0.5\textwidth]{Solution/AMC_10_10A_2023_22_Solution_Figure_1.png}
                    \end{figure}
                    

 Let $O$ be the center of the midpoint of the line segment connecting both the centers, say $A$ and $B$.

 Let the point of tangency with the inscribed circle and the right larger circles be $T$. 

 Then $OT = BO + BT = BO + AT - \frac{1}{2} = \frac{1}{4} + 1 - \frac{1}{2} = \frac{3}{4}.$

 Since $C_4$ is internally tangent to $C_1$, center of $C_4$, $C_1$ and their tangent point must be on the same line.

 Now, if we connect centers of $C_4$, $C_3$ and $C_1$/$C_2$, we get a right angled triangle.

 Let the radius of $C_4$ equal $r$. With the pythagorean theorem on our triangle, we have

 
 \[\left(r+\frac{3}{4}\right)^2+\left(\frac{1}{4}\right)^2=(1-r)^2\]

 Solving this equation gives us

 
 \[r = \boxed{\textbf{(D) } \frac{3}{28}}\]




\textbf{Solution 2: Estimation}

Draw 2 lines from the middle segment's endpoints to an intersection of $C_1$ and $C_2$. Draw the perpendicular bisector of the middle segment until it reaches one of the intersections. We see this becomes a right triangle with a hypotenuse of $1$ and one leg length of $\dfrac14$. Using the Pythagorean theorem, the other leg length is $\dfrac{\sqrt{15}}{4}$. Also drawing $C_1$'s radius along the segment of $\dfrac12$ allows us to find the that radius of the middle circle is $\dfrac34$. 

 Finally notice that $r_{C_3} + 2 r_{C_4}$ is really close to $\dfrac{\sqrt{15}}{4}$. $2 r_{C_4} = \dfrac{\sqrt{15} - 3}{4}$, so $r_{C_4} = \dfrac{\sqrt{15} - 3}{8}$. This is really close to $\dfrac18$ and the closest answer to that is $\boxed{\text{D) }\dfrac{3}{28}}$



	\item (\textit{AMC 10 10A 2023}) If the positive integer $c$ has positive integer divisors $a$ and $b$ with $c = ab$, then $a$ and $b$ are said to be $\textit{complementary}$ divisors of $c$. Suppose that $N$ is a positive integer that has one complementary pair of divisors that differ by $20$ and another pair of complementary divisors that differ by $23$. What is the sum of the digits of $N$?


 $\textbf{(A) } 9 \qquad \textbf{(B) } 13\qquad \textbf{(C) } 15 \qquad \textbf{(D) } 17 \qquad \textbf{(E) } 19$

 \textbf{Answer: }C



\textbf{Solution 1}

Consider positive integers $a, b$ with a difference of $20$. Suppose $b = a-20$. Then, we have $(a)(a-20) = n$. If there is another pair of two integers that multiply to $n$ but have a difference of 23, one integer must be greater than $a$, and the other must be smaller than $a-20$. We can create two cases and set both equal. 

 We have $(a)(a-20) = (a+1)(a-22) \text{ or } (a+2)(a-21).$ Note that if we go further to $(a+3)(a-20)$ and beyond, that would violate the condition that one of the two integers must be smaller than $a-20.$

 Starting with the first case, we have $a^2-20a = a^2-21a-22$, or $0=-a-22$, which gives $a=-22$, which is not possible. The other case is $a^2-20a = a^2-19a-42$, so $a=42$. Thus, our product is $(42)(22) = (44)(21)$, so $n = 924$. Adding the digits, we have $9+2+4 = \boxed{\textbf{(C) } 15}$.

 -Sepehr2010, mathboy282, the\_eaglercraft\_grinder

 



\textbf{Solution 2}

We have 4 integers in our problem. Let's call the smallest of them $a$. $a(a+23) =$ either $(a+1)(a+21)$ or $(a+2)(a+22)$. So, we have the following:

 $a^2 + 23a = a^2 + 22a +21$ or 

 $a^2+23a = a^2 + 24a +44$.

 The second equation has negative solutions, so we discard it. The first equation has $a = 21$, and so $a + 23 = 44$. If we check $(a+1)(a+21)$ we get $22 \cdot 42 = 21 \cdot 44$. $44$ is $2$ times $22$, and $42$ is $2$ times $21$, so our solution checks out. Multiplying $21$ by $44$, we get $924 \implies 9 + 2 + 4 = \boxed{\textbf{(C) 15}}$. 

 



\textbf{Solution 3 (less words)}

From the problems, it follows that

 
 \begin{align*} x(x+20)&=y(y+23) = N\\ x^2+20x&=y^2+23y\\ 4x^2+4\cdot20x &= 4y^2+4\cdot23y\\ 4x^2+4\cdot20x+20^2-20^2 &= 4y^2+4\cdot23y+23^2-23^2\\ (2x+20)^2-20^2 &= (2y+23)^2-23^2\\ 23^2-20^2 &= (2y+23)^2-(2x+20)^2\\ (23+20)(23-20) &= (2y+23+2x+20)(2y+23-2x-20)\\ 43\cdot 3 &= (2y+2x+43)(2y-2x+3)\\ 129\cdot 1 &= (2y+2x+43)(2y-2x+3)\\ \end{align*}
Since both $(2y+2x+43)$ and $(2y-2x+3)$ must be integer, we get two equations.
 \begin{align} 129 \text{ or } 43 &= (2y+2x+43)\\ 1 \text{ or } 3 &= 2y-2x+3\\ \end{align}43 & 3 yields (0,0) which is not what we want.129 & 1 yields (22,21) which is more interesting.

 Simplifying the equations, we get:
 \begin{align*} x+y &= 43\\ x-y &= 1\\ x=22&, y=21\\ N &= (22)(22+20) = 924. \end{align*}


 So, the answer is $\boxed{\textbf{(C) 15}}$.

 


 



\textbf{Solution 4}

Say one factorization is $n(n+23).$ The two cases for the other factorization are $(n+1)(n+21)$ and $(n+2)(n+22).$ We know it must be the first because of AM-GM intuition: lesser factors of a number are closer together than larger factors of a number. (We can also try both and see which works.) Thus, $n(n+23)=(n+1)(n+21)$ and we find that $n=21,N=924$ meaning the answer is $\boxed{\textbf{(C) }15}.$




\textbf{Solution 5}

Since we are given that some pairs of divisors differ by 20 and 23 and we can let the pair be $(x-10)$ and $(x+10)$ as well as $(y-\frac{23}{2})$ and $(y+\frac{23}{2})$ . We also know the product of both the complementary divisors give the same number so $(x-10)(x+10)=(y-\frac{23}{2})(y+\frac{23}{2})$ . Now we let $y=\frac{a}{2}$. Then we substitute and get $x^2-100=\frac{(a^2-529)}{4}$. Finally we multiply by 4 and get $4x^2-a^2=-129, a^2-4x^2=129$. Then we use differences of squares and get $a$+$2x$=129, $a$-$2x$=1. We finish by getting $a=$65 and $x=32$. So $(42)(22) = 924$ Adding the digits, we have $9+2+4 = \boxed{\textbf{(C) } 15}$.

 
Nunber sense note: To avoid tedious multiplication of 2-digit numbers, observe that $n = (42)(22) = (6)(7)(2)(11)$, and $(6)(7)(2) = 84$, and the sum of the digits of $11$ is $2$,  so the sum of the digits of $n$ is equivalent to $(8+4)(2) \equiv 24 \equiv 15 \pmod 9$.  The only equivalent answer choice is  $\boxed{15}$. ~oinava

 



\textbf{Solution 6}

$N$ can be written $N = \left( a - 10 \right) \left( a + 10 \right)$ with a positive integer $a > 10$ and $N = \left( \frac{2b + 1}{2} - \frac{23}{2} \right) \left( \frac{2b + 1}{2} + \frac{23}{2} \right)$ with a positive integer $b > 11$.

 The above equations can be reorganized as
 \[ \left( 2b + 1 + 2 a \right) \left( 2 b + 1 - 2 a \right) = 43 \cdot 3 . \]

 The only solution is $2b + 1 + 2a = 129$ and $2b + 1 - 2a = 1$.Thus, $a = b = 32$.Therefore, $N = 924$.So the sum of the digits of $N$ is $9 + 2 + 4 = \boxed{\textbf{(C)}~15}$.

 



\textbf{Solution 7}

We can write $N$ as $a(a+20)$ or $b(b+23)$ where $a$ and $b$ are divisors of $N.$ Since $a(a+20) = b(b+23),$ we know that $a^2 + 20a - b^2 - 23b = 0$, and we can view this as a quadratic in $a.$ 

 Since the solution for $a$ must be an integer, the discriminant for this quadratic must be a perfect square and therefore $20^2 - 4(-b^2 - 23b) = (2c)^2 = 400 + 4b^2 + 92b$ so $b^2 + 23b -c^2 + 100 = 0.$ 

 Since the discriminant of this quadratic in $b$ must also be a perfect square we know that $23^2 - 4(-c^2+100) = d^2$ which we can simplify as $d^2 - 4c^2 = (d-2c)(d+2c) = 129.$ Since they are both positive integers $d - 2c$ and $d + 2c$ are factors of $129 = 3 \cdot 43$ so $d - 2c = 1$ and $d + 2c = 129$ or $d - 2c = 3$ and $d - 2c = 43.$ 

 These systems of equations give us $(c,d) = (32,65)$ and $(c,d) = (10,23)$ respectively, if we plug our values for $c$ into the equation for $b$ we get $b^2 + 23b - 924 = 0$ and $b^2 + 23b = 0$ respectively. The first equation gives us $b = 21$ or $b = -44$ and the second gives us $b = 0$ or $b = -23$, since $b$ is positive we know that $b = 21$ and $N = (21)(21 + 23) = 924$, therefore the sum of the digits of $N$ is $9 + 2 + 4 = \boxed{\textbf{(C) 15}}.$




\textbf{Solution 8 (Trial and Error)}

Consider the numbers of the form $a(a+20)$. Since $b(b+23)$ is always even, $a$ is even. Thus, for $a \ge 2$, we calculate $a(a+20)$ for even values of $a$. Then, we check if it can also be represented as a product of numbers that differ by $23$. Checking, we see that $22 \cdot 42 = 21 \cdot 44 = 924$ works. Thus, the answer is $9 + 2 + 4 = \boxed{\textbf{(C) 15}}$




\textbf{Solution 9}

$n(n+20)=m(m+23) \Longrightarrow n^2+20n=m^2+23m \Longrightarrow n^2-m^2+20n-20m=3m$. Factoring, $(n+m)(n-m)+20(n-m)=3m \Longrightarrow (n-m)(n+m+20)=3m$. Let $n-m=a>0$ because clearly $n>m$. Then. $a(2m+20+a)=3m$. Note that since $20+a>0$, if $a\geq2$, then the equation is $4m+a(20+a)>3m$, so $a=1$. Plugging this back, we get $2m+21=3m \Longrightarrow m=21$ and $n=22$. Now we find $N$ as $22*42=924$ so the answer is $15$.

 -Magnetoninja

 



\textbf{Solution 10 (Elementary Algebra)}

We are given $a(a+20) = b(b+23) = N$, and we are trying to find $N$.Simplifying the equation yields

 
 \[a^2+20a=b^2+23b.\]

 Completing the square for $a$, we get

 
 \[(a+10)^2=b^2+23b+100.\]

 Now, we will use the squeezing method to find $b$. We can set a bound on $b^2+23b+100$: 

 
 \[b^2+20b+100<b^2+23b+100<b^2+24b+144.\]

 because the expression is bigger than the lower bound by $3b$ and less than the upper bound by $b+44$.Simplifying gets:

 
 \[(b+10)^2<b^2+23b+100<(b+12)^2.\]

 Since we know that the expression is a perfect square, it must be equal to the only square in between $b+10$ and $b+12$, i.e. $b+11$:

 
 \[(b+11)^2=b^2+23b+100.\]
 \[b^2+22b+121=b^2+23b+100\]
 \[b=21\]Therefore $N=21(21+23)=21(44)=924.$ The sum of the digits is $9+2+4=\boxed{\textbf{(C) 15}}$




\textbf{Solution 11 (Algebra, and some hopeful guesswork)}

Let $N=a(a-20)=b(b-23)$.We get $a^2 - 20a = b^2 - 23b$. Completing the square gives $(a-10)^2 - 100 = (b - \frac{23}{2})^2 - \frac{529}{4}$. We can add $\frac{529}{4}$ and then multiply by $4$ to get $4(a-10)^2 + 129 = 4(b - \frac{23}{2})^2 \rightarrow (2a-20)^2 + 129 = (2b - 23)^2$. We can rewrite $129$ as $2(64) + 1$, which makes the LHS look similar to the form $x^2 + 2x + 1$, where $x = 64$. So we guess that $(2a-20) = 64$, and $(2b-23) = 64 + 1 = 65$. Solving these two equations gives $a=42$ and $b=44$. This means $N=42(22)=(44)(21)=924$, so the sum of the digits of $N$ is $9+2+4=\boxed{\textbf{(C) 15}}$, which is our final answer.

 


 


	\item (\textit{AMC 10 10A 2023}) Six regular hexagonal blocks of side length $1$ unit are arranged inside a regular hexagonal frame. Each block lies along an inside edge of the frame and is aligned with two other blocks, as shown in the figure below. The distance from any corner of the frame to the nearest vertex of a block is $\frac{3}{7}$ unit. What is the area of the region inside the frame not occupied by the blocks?


 
                    \begin{figure}[H]
                        \centering
                        \includegraphics[width=0.5\textwidth]{Problem/AMC_10_10A_2023_24_Problem_Figure_1.png}
                    \end{figure}
                    
$\textbf{(A)}~\frac{13 \sqrt{3}}{3}\qquad\textbf{(B)}~\frac{216 \sqrt{3}}{49}\qquad\textbf{(C)}~\frac{9 \sqrt{3}}{2} \qquad\textbf{(D)}~ \frac{14 \sqrt{3}}{3}\qquad\textbf{(E)}~\frac{243 \sqrt{3}}{49}$

 \textbf{Answer: }C



\textbf{Solution 1}


                    \begin{figure}[H]
                        \centering
                        \includegraphics[width=0.5\textwidth]{Solution/AMC_10_10A_2023_24_Solution_Figure_1.png}
                    \end{figure}
                    

 
Examining the red isosceles trapezoid with $1$ and $\dfrac{3}{7}$ as two bases, we know that the side lengths are $\dfrac{4}{7}$ from $30-60-90$ triangle.

 We can conclude that the big hexagon has side length 3.

 Thus the target area is: area of the big hexagon - 6 * area of the small hexagon.$\dfrac{3\sqrt{3}}{2}(3^2-6\cdot1^2) = \dfrac{3\sqrt{3}}{2}(3) = \boxed{\textbf{(C)}~\frac{9 \sqrt{3}}{2}}$




\textbf{Solution 1.1}

We can extend the line of the parallelogram to the end until it touches the next hexagon and it will make a small equilateral triangle and a longer parallelogram. We can prove that one side of the tiny equilateral triangle is 4/7 by playing around with angles and the parallelogram because it is parallel, we can then use the whole side of the hexagon which is one and subtract 3/7 which is one side of the equilateral triangle which is 4/7. That means the whole side of the big hexagon length is 3 and we can continue with solution 1.

 
                    \begin{figure}[H]
                        \centering
                        \includegraphics[width=0.5\textwidth]{Solution/AMC_10_10A_2023_24_Solution_Figure_2.png}
                    \end{figure}
                    -BrandonChen

 



\textbf{Solution 1.3}

An alternative diagram, perhaps slightly more intuitive, can be drawn by extending segments to form two equilateral triangles (shown in blue) with side length 1.
                    \begin{figure}[H]
                        \centering
                        \includegraphics[width=0.5\textwidth]{Solution/AMC_10_10A_2023_24_Solution_Figure_3.png}
                    \end{figure}
                    

 The red segment is given to have length \(3/7\). Furthermore, we can find yet another equilateral triangle within the bottom blue triangle that has side length \(1-3/7 = 4/7\), shown in dashed magenta. Adding all four segments (red, blue, blue, magenta), we find the side length of the larger hexagon to be \(3/7+1+1+4/7 = 3\). Then we can continue with Solution 1.

 (diagram borrowed mostly from ~Technodoggo's Solution 1)

 



\textbf{Solution 2 (Not rigorous)}

Note that one can "slide' the small hexagons along their respective edges, and either by sliding them to the center or to the corners, and thus getting that the side length of the larger hexagon is 3. The rest proceeds the same as solution 1.

 



\textbf{Solution 2.1 (Clarification)}

Notice that when sliding the smaller hexagon along the edge, we see that the contact edge withe the smaller hexagon "in front" of it is $60^{\circ}$, thus meaning the hexagon "in front" is pushed at a speed $\sin{60^{\circ}}$ times the actual speed of the hexagon. We can preform a similar analysis on the hexagon that is being pushed and get that the speed at which that hexagon is moving is $\frac{1}{\sin{60^{\circ}}}$ times the speed it pushed by. As we can see, the two factors cancel out and by the same argument, every small hexagon can move at the same speed while mantaining an edge of contact with the two adjacent hexagons.

 
                    \begin{figure}[H]
                        \centering
                        \includegraphics[width=0.5\textwidth]{Solution/AMC_10_10A_2023_24_Solution_Figure_4.png}
                    \end{figure}
                    
                    \begin{figure}[H]
                        \centering
                        \includegraphics[width=0.5\textwidth]{Solution/AMC_10_10A_2023_24_Solution_Figure_5.png}
                    \end{figure}
                    \href{https://artofproblemsolving.com/community/c5h3194971p29157634}{Diagram by ihatemath123}

 \textbf{Note} The number $\frac{3}{7}$ is irrelevant to solve the problem.In fact, if the smaller hexagons have side length $x$, the side length of the large hexagon will always be $3x$.

 



\textbf{Solution 2}

We put the diagram to a complex plane, with the center of the outside hexagon at the origin.We denote by $s$ the length of each side of the outside hexagon.

 The complex number of the upper left vertex of the upper right small hexagon is
 \begin{align*} s e^{i 30^\circ} + \frac{3}{7} e^{i \left( 90^\circ + 60^\circ \right)} - \sqrt{3} . \end{align*}


 The complex number of the upper right vertex of the top small hexagon is
 \begin{align*} s i + \frac{3}{7} e^{i \left( \pi + \frac{\pi}{6} \right)} + e^{- i \frac{\pi}{6}} . \end{align*}


 The above two vertices are on the same vertical line. So their real part values are the same.By solving this equation, we get $s = 3$.

 Therefore, the area of the region not occupied by the blocks is
 \begin{align*} 6 \cdot \frac{\sqrt{3}}{4} s^2 - 6 \cdot 6 \cdot \frac{\sqrt{3}}{4} 1^2 = \boxed{\textbf{(C) } \frac{9 \sqrt{3}}{2}} . \end{align*}





\textbf{Solution 3 (no words)}


                    \begin{figure}[H]
                        \centering
                        \includegraphics[width=0.5\textwidth]{Solution/AMC_10_10A_2023_24_Solution_Figure_6.png}
                    \end{figure}
                    
$(3^2-6\times 1^2)\times\frac{3\sqrt{3}}{2}=(9-6\times 1)\times\frac{3\sqrt{3}}{2}=(9-6)\times\frac{3\sqrt{3}}{2}=3\times\frac{3\sqrt{3}}{2}=\boxed{\textbf{(C) } \frac{9 \sqrt{3}}{2}}$




\textbf{Solution 4 (educated guessing)}

Since the area of each hexagon is $\frac{3 \sqrt{3}}{2}$, the answer is just the area of the large hexagon minus $9\sqrt{3}$. The side length of the large hexagon should be an integer or fraction since we are given integers and fractions, and solving this problem shouldn't require any square roots. Therefore, the correct answer plus $9\sqrt{3}$ should be the product of $\frac{3 \sqrt{3}}{2}$ and a square number. Only choice $\text{C}$ satisfies this, since $\frac{9 \sqrt{3}}{2} + 9\sqrt{3}=\frac{27 \sqrt{3}}{2}=3^2\times\frac{3 \sqrt{3}}{2}$. Alternatively, we can also deduce $\text{C}$ is correct since it is the only fraction whose denominator is ${2}$. Therefore, the answer is (probably) $\boxed{\textbf{(C)}~\frac{9 \sqrt{3}}{2}}$

 -faliure167

 


	\item (\textit{AMC 10 10A 2023}) If $A$ and $B$ are vertices of a polyhedron, define the distance $d(A, B)$ to be the minimum number of edges of the polyhedron one must traverse in order to connect $A$ and $B$. For example, $\overline{AB}$ is an edge of the polyhedron, then $d(A, B) = 1$, but if $\overline{AC}$ and $\overline{CB}$ are edges and $\overline{AB}$ is not an edge, then $d(A, B) = 2$. Let $Q$, $R$, and $S$ be randomly chosen distinct vertices of a regular icosahedron (regular polyhedron made up of $20$ equilateral triangles). What is the probability that $d(Q, R) > d(R, S)$?


 $\textbf{(A) }\frac{7}{22}\qquad\textbf{(B) }\frac{1}{3}\qquad\textbf{(C) }\frac{3}{8}\qquad\textbf{(D) }\frac{5}{12}\qquad\textbf{(E) }\frac{1}{2}$

 \textbf{Answer: }A



\textbf{Solution 1}

To find the total amount of vertices we first find the amount of edges, and that is $\frac{20 \times 3}{2}$. Next, to find the amount of vertices we can use Euler's characteristic, $V - E + F = 2$, and therefore the amount of vertices is $12$ 

 So there are $\binom{12}{3} = 220$ ways to choose 3 distinct points.

 Now, the furthest distance we can get from one point to another point in an icosahedron is 3. Which gives us a range of $1 \leq d(Q, R), d(R, S) \leq 3$

 With some case work, we get two cases:

 Case 1: $d(Q, R) = 3; d(R, S) = 1, 2$

 Since we have only one way to choose Q, that is, the opposite point from R, we have one option for Q and any of the other points could work for S.

 Then, we get $12 \times 1 \times 10 = 120$ (ways to choose R × ways to choose Q × ways to choose S)

 Case 2: $d(Q, R) = 2; d(R, S) = 1$

 We can visualize the icosahedron as 4 rows, first row with 1 vertex, second row with 5 vertices, third row with 5 vertices and fourth row with 1 vertex. We set R as the one vertex on the first row, and we have 12 options for R. Then, Q can be any of the 5 points on the third row and finally S can be one of the 5 points on the second row.

 Therefore, we have $12 \times 5 \times 5 = 300$ (ways to choose R × ways to choose Q × ways to choose S)

 In total, we have $420$ ways, but we must divide by $3!$ to account for permutations, giving us $70/220 = \boxed{\textbf{(A) } 7/22}.$




\textbf{Solution 2 (Cheese + Actual way)}

In total, there are $\binom{12}{3}=220$ ways to select the points. However, if we look at the denominators of $B,C,D$, they are $3,8,12$ which are not divisors of $220$. Also $\frac{1}{2}$ is impossible as cases like $d(Q, R) = d(R, S)$ exist. The only answer choice left is $\boxed{\textbf{(A) } \frac{7}{22}}$

 Note: this cheese is actually wrong because the total number of ways to select the points is actually $12 \times 11 \times 10 = 1320$ as order matters, so all denominators are possible. Rather, you can arrive at the same conclusion by fixing R WLOG, leading to $11 \times 10 = 110$ ways in total, which works for the original cheese.. ~awesomeguy856

 (Actual way)

 Fix an arbitrary point, to select the rest $2$ points, there are $11\cdot 10 =110$ ways. To make $d(Q, R)=d(R, S), d=1/2$. Which means there are in total $2\cdot 5 \cdot 4 = 40$ ways to make the distance the same. $\frac{1}{2}(1-\frac{40}{110})= \boxed{\textbf{(A) } \frac{7}{22}}$~bluesoul

 



\textbf{Solution 3}

We can imagine the icosahedron as having 4 layers. 1 vertex at the top, 5 vertices below connected to the top vertex, 5 vertices below that which are 2 edges away from the top vertex, and one vertex at the bottom that is 3 edges away. WLOG because the icosahedron is symmetric around all vertices, we can say that R is the vertex at the top. So now, we just need to find the probability that S is on a layer closer to the top than Q. We can do casework on the layer S is on to get 
 \[\frac{5}{11} \cdot \frac{6}{10} + \frac{5}{11} \cdot \frac{1}{10} = \frac{35}{110} = \frac{7}{22}\]So the answer is $\boxed{\textbf{(A) }\frac{7}{22}}$. 

 -awesomeparrot

 



\textbf{Solution 4}

We can actually see that the probability that $d(Q, R) > d(R, S)$ is the exact same as $d(Q, R) < d(R, S)$ because $d(Q, R)$ and $d(R, S)$ have no difference. (In other words, we can just swap Q and S, meaning that can be called the same probability-wise.) Therefore, we want to find the probability that $d(Q, R) = d(R, S)$.

 WLOG, we can rotate the icosahedron so that R is the top of the icosahedron. Then we can divide this into 2 cases:

 1. They are on the second layer

 There are 5 ways to put one point, and 4 ways to put the other point such that $d(Q, R) = d(R, S) = 1$. So, there are $5 \cdot 4 = 20$ ways to put them on the second layer.

 2. They are on the third layer

 There are 5 ways to put one point, and 4 ways to put the other point such that $d(Q, R) = d(R, S) = 2$. So, there are $5 \cdot 4 = 20$ ways to put them on the third layer.

 The total number of ways to choose P and S are $11 \cdot 10 = 110$ (because there are 12 vertices), so the probability that $d(Q, R) = d(R, S)$ is $\frac{20+20}{110} = \frac{4}{11}$.

 Therefore, the probability that $d(Q, R) > d(R, S)$ is $\frac{1 - \frac{4}{11}}{2} = \boxed{\textbf{(A) }\frac{7}{22}}$




\textbf{Solution 5}

We know that there are $20$ faces. Each of those faces has $3$ borders (since each is a triangle), and each edge is used as a border twice (for each face on either side). Thus, there are $\dfrac{20\cdot3}2=30$ edges. 

 Euler's formula states that $V-E+F=2$ for all convex polyhedra, so we know that $2-F+E=12$ vertices. 

 The answer can be counted by first counting the number of possible paths that will yield $d(Q, R) > d(R, S)$ and dividing it by $12\cdot11\cdot10$ (or $\dbinom{12}3$, depending on the approach). In either case, one will divide by $11$ somewhere in the denominator. We can then hope that there will be no factor of $11$ in the numerator (which would cancel the $11$ in the denominator out), and answer the only option that has an $11$ in the denominator: $\boxed{\textbf{(A) }\frac{7}{22}}$. 

 Additional note by "Fruitz": One can eliminate $1/2$ by symmetry if you swap the new sign.

 Another note by "andliu766": A shorter way to find the number of vertices and edges is to use the fact that the MAA logo is an icosahedron. :)

 



\textbf{Solution 6 (Case Work(similar to 3))}


                    \begin{figure}[H]
                        \centering
                        \includegraphics[width=0.5\textwidth]{Solution/AMC_10_10A_2023_25_Solution_Figure_1.png}
                    \end{figure}
                    WLOG, let R be at the top-most vertex of the icosahedron. There are $2$ cases where $d(Q, R) > d(R, S)$.

 Case 1: $Q$ is at the bottom-most vertex

 If $Q$ is at the bottom-most vertex, no matter where $S$ is, $d(Q, R) > d(R, S)$. The probability that $Q$ is at the bottom-most vertex is $\frac{1}{11}$

 Case 2: $Q$ is at the second layer

 If $Q$ is at the second layer, $S$ must be at the first layer, for $d(Q, R) > d(R, S)$ to be true. The probability that $Q$ is at the second layer, and $S$ is at the first layer is $\frac{5}{11} \cdot \frac{5}{10} = \frac{5}{22}$

 \[\frac{1}{11} + \frac{5}{22} = \boxed{\textbf{(A) }\frac{7}{22}}\]

 \textbf{Supplement} First, we have 20 faces of equilateral triangles. But each edge is shared by two faces hence there are thus $\frac{20 \cdot 3}{2}=30$ edges. 

 Now by Euler's formula $V-E+F=2 \Rightarrow V-30+20=2 \Rightarrow V=12.$ 

 Furthermore, the number of triangles ($n$) that meet at one vertex can be found as follows. Note that there are $12$ total vertices. Without accounting for overcounting, $\text{total number of vertices} = \text{number of faces} \cdot \text{vertices per face} = 20 \cdot 3 = 60.$ But since we now know that there are only $12$ vertices and that each vertex is shared by $n$ triangles, you find that there are 5 triangles that meet at one vertex.

 Once we pick R as the topmost vertex, there are 11 possible vertices remaining for $Q.$ Continue as follows in this solution.

 



\textbf{Solution 7 (efficient)}

Since the icosahedron is symmetric polyhedron, we can rotate it so that R is on the topmost vertex. Since Q and S basically the same, we can first count the probability that $d(Q,R) = d(R,S)$.

 $\mathfrak{Case} \ \mathfrak{1}: d(Q,R) = d(R,S) = 1$

 There are 5 points $P$ such that $d(Q,P) = 1$. There is $5 \times 4 = \boxed{20}$ ways to choose Q and S in this case.

 $\mathfrak{Case} \ \mathfrak{2}: d(Q,R) = d(R,S) = 2$

 There are 5 points $P$ such that $d(Q,P) = 2$. There is $5 \times 4 = \boxed{20}$ ways to choose Q and S in this case.

 $\mathfrak{Case} \ \mathfrak{3}: d(Q,R) = d(R,S) = 3$

 There is 1 point $P$ such that $d(Q,P) = 3$. There is $1 \times 0 = \boxed{0}$ ways to choose Q and S in this case.

 $\mathfrak{Final} \ \mathfrak{solution}$

 There are 11 points $P$ that are distinct from R. There is $11 \times 10 = \boxed{110}$ ways to choose Q and S. There is $20 + 20 + 0 = \boxed{40}$ ways to choose Q and S such that $d(Q,R) = d(R,S)$. There is $\frac{110-40}{2} = 35$ ways to choose Q and S such that $d(Q, R) > d(R, S)$. The probability that $d(Q, R) > d(R, S)$ is therefore $\frac{35}{110} = \frac{7}{22}$ which corresponds to answer choice $\boxed{A}$




\textbf{Solution 8 (Dodecahedron Schlegel (Planar Graph) Diagram)}


                    \begin{figure}[H]
                        \centering
                        \includegraphics[width=0.5\textwidth]{Solution/AMC_10_10A_2023_25_Solution_Figure_2.png}
                    \end{figure}
                    

 Instead of thinking this problem as vertices of a icosahedron, think of it as faces of a dodecahedron. Our middle face is $R$ and we first choose $Q$ and $S$ to be the faces with distance $1$ from the middle face. There are $\binom{5}{2}$ ways of doing that. Then we choose $Q$ and $S$ to be the faces with distance $2$ from the middle face. There are $\binom{5}{2}$ ways of doing that. Our answer is therefore $\frac{1}{2}\left(1-\frac{2\binom{5}{2}}{\binom{11}{2}}\right)=\boxed{\textbf{(A) }\frac{7}{22}}$

 
Note: You can do the same analysis on the vertices of the icosahedron using its diagram, which is less to draw because it has fewer (12) vertices and the same number (30) of edges. But this diagram is either less symmetrical or non-planar. (There's no good place to put the last vertex (diametrically opposite the first/central vertex), except in the same place as the first vertex.)

 



\textbf{Solution 9 (Symmetry)}

Note that the vertices of an icosahedron are in a $1-5-5-1$ configuration.

 We will denote the probability of event $A$ happening as $P(A)$

 We should also note that $P(d(Q, R) > d(R, S)) + P(d(Q, R) < d(R, S)) + P(d(Q, R) = d(R, S)) = 1$

 Also, $P(d(Q, R) > d(R, S)) = P(d(Q, R) < d(R, S))$ since $Q$ and $S$ are basically the same.

 Therefore, what we are trying to find is just $\frac{1-P(d(Q, R) = d(R, S))}{2}$ which can easily be computed when we fix $R$ at the top vertex in the $1-5-5-1$ structure.

 Since all the points are distinct (stated in problem), the conditions are only satisfied where $Q$ and $S$ are at $5$s of the $1-5-5-1$ structure.

 $P(d(Q, R) = d(R, S)) = 2 \cdot \frac{5}{11} \cdot \frac{4}{10} = \frac{4}{11}$ since we have to choose one of the $5$ out of $11$ possible vertices and another one of $4$ out of $10$ two times.

 And the answer is $\frac{1-\frac{4}{11}}{2} = \boxed{ \textbf{(A) } \frac{7}{22}}$



\end{enumerate}